\documentclass{ieeeaccess}
\usepackage{cite}
\usepackage{amsmath,amssymb,amsfonts}
\usepackage{algorithmic}
\usepackage{graphicx}
\usepackage{textcomp}
\usepackage{caption}
\usepackage{mathptmx}
\usepackage{hyperref}
\usepackage[english]{babel}
\usepackage{microtype} 
\usepackage{xurl}
\usepackage{booktabs}
\usepackage{tabularx}
\usepackage{placeins}
\usepackage{bm}
\usepackage{bbm}
\usepackage{ragged2e}
\makeatletter
\AtBeginDocument{\DeclareMathVersion{bold}
\SetSymbolFont{operators}{bold}{T1}{times}{b}{n}
\SetSymbolFont{NewLetters}{bold}{T1}{times}{b}{it}
\SetMathAlphabet{\mathrm}{bold}{T1}{times}{b}{n}
\SetMathAlphabet{\mathit}{bold}{T1}{times}{b}{it}
\SetMathAlphabet{\mathbf}{bold}{T1}{times}{b}{n}
\SetMathAlphabet{\mathtt}{bold}{OT1}{pcr}{b}{n}
\SetSymbolFont{symbols}{bold}{OMS}{cmsy}{b}{n}
\renewcommand\boldmath{\@nomath\boldmath\mathversion{bold}}}
\makeatother

\def\BibTeX{{\rm B\kern-.05em{\sc i\kern-.025em b}\kern-.08em
    T\kern-.1667em\lower.7ex\hbox{E}\kern-.125emX}}

%Your document starts from here ___________________________________________________
\begin{document}
\history{Date of publication xxxx 00, 0000, date of current version xxxx 00, 0000.}
\doi{10.1109/ACCESS.2024.0429000}


\title{On Reliable Detection of AI-Generated Multimedia Content}
\author{\uppercase{Ansuman Nayak}\authorrefmark{1}, and 
\uppercase{Vijay K. Madisetti}\authorrefmark{2}, \IEEEmembership{Fellow, IEEE}}

\address[1]{College of Computing, Georgia Institute of Technology, Atlanta, GA 30332 USA (e-mail: anayak306@gatech.edu)}
\address[2]{School of Cybersecurity and Privacy, Georgia Institute of Technology, Atlanta, GA 30332 USA (e-mail: vkm@gatech.edu)}

\markboth
{Nayak \headeretal: On Reliable Detection of AI-Generated Multimedia Content}
{Nayak \headeretal: On Reliable Detection of AI-Generated Multimedia Content}

\corresp{Corresponding author: Dr. Vijay Madisetti (e-mail: vkm@gatech.edu).}


\begin{abstract}
With the rapid proliferation of AI-generated multimedia content, the growing need for reliable detection of AI-generated media, including deepfakes and misinformation, has been largely unmet. Uniformly applying large vision-language models (VLMs) to all content can be prohibitively expensive. In this work, we formulate detection as a model selection problem over a set of complementary detectors. We show that disagreement between artifact-based and semantic detectors provides a strong, parameter-free indicator of input difficulty, and that this signal remains stable even when the semantic probe faces deployment shift to unseen generator families. We construct a selective-computation cascade in which lightweight detectors process all inputs and a stronger verifier is invoked only for disagreement cases. On in-domain benchmarks, a well-trained semantic detector achieves the best cost-accuracy tradeoff without VLM escalation. The cascade's utility emerges under generator evolution, but only when the verifier remains stronger than the stale semantic detector on the routed disagreement subset. On SDXL and FLUX, escalating to GPT-5.5 only on disagreement cases recovers up to 7.9 percentage points of accuracy over a stale semantic probe at 42--51\% of the cost of escalating on every input, closing most but not all of the gap to a uniform always-escalate policy (which reaches 97--100\% accuracy at full cost). In contrast, under a more distant BigGAN shift, escalation provides no benefit, showing that routing effectiveness depends on verifier strength and distributional distance rather than disagreement alone. We additionally show that BigGAN's near-ceiling in-domain artifact-detector accuracy is substantially attributable to a native-resolution shortcut in the benchmark rather than generative-artifact detection, identified via a resolution-only control classifier and a resolution-matching ablation. We further introduce a verifier validation gate with a prespecified, uncertainty-aware decision rule. Evaluated on a genuinely held-out calibration/test split (not the full-subset oracle used in earlier validation), the gate reaches at least 90\% test-set agreement with the oracle decision in 21 of 32 eligible configurations at $N=100$ and 26 of 32 at $N=200$, and reduces harmful escalations by an order of magnitude relative to a point-estimate decision rule; a three-way detector selector is shown to be equivalent to the binary gate as a decision procedure.

\end{abstract}

\begin{keywords}
AI-generated image detection, deepfake 
detection, multi-model disagreement, distribution shift, 
vision-language models, forensic routing, CLIP embeddings
\end{keywords}

\titlepgskip=-21pt
\maketitle

% Sections
\section{Introduction}

The rapid advancement of generative models for images and videos has significantly increased the realism and accessibility of synthetic media. Diffusion models~\cite{Wang_2023_ICCV}, GAN-based generators~\cite{wang2020cnn}, and multimodal large language models (MLLMs) can now generate highly convincing content with minimal effort. Although this opens up useful applications, it also raises concerns about misinformation, impersonation, and fraud. As a result, reliably detecting AI-generated media has become an important problem that still needs to be addressed.

One of the main challenges is generalization. Detection models often perform well on the generator families they were trained on, but accuracy degrades significantly when applied to new or unseen generators~\cite{NEURIPS2023_f4d4a021}. Artifact-based detectors exploit low-level frequency and synthesis patterns introduced during generation, while semantic models capture higher-level structural representations. These two signal types fail in different ways, suggesting that combining them provides more robust coverage than either approach alone.


In practical deployment settings, compute cost is a significant concern. High-capacity
vision-language models (VLMs) provide strong reasoning capabilities but are expensive to apply at scale~\cite{chen2024frugalgpt}. Applying these models uniformly to every input is often infeasible in real-time or large-scale systems. Existing approaches typically apply a
fixed detector or ensemble without adapting computation to input difficulty. The central challenge is identifying which inputs require expensive verification rather than applying a
uniform computation to all inputs.


In this work, we address this challenge by modeling AI-generated media detection as a model selection problem over complementary detectors. We find that the disagreement between an artifact-based detector and a semantic detector serves as a proxy for input difficulty and requires no calibration. When the detectors agree, the predictions are highly reliable, whereas disagreements consistently identify cases where cheap detectors are insufficient.

Unlike prior approaches that rely on calibrated confidence scores or learned routing policies, disagreement provides a parameter-free signal that remains reliable under distribution shift. We design a cascade that routes inputs based on a signal where lightweight detectors process all inputs, and a stronger verifier is invoked only on disagreement cases. When the semantic probe is trained in-domain, it already achieves strong performance without VLM escalation. The cascade's primary value is under deployment shift, where it recovers
accuracy over a stale semantic probe while avoiding the majority of verifier invocations. This reframes the detection problem as identifying which inputs require expensive computation, rather than uniformly increasing model complexity.

\subsection{Research Questions}

This work is guided by the following research questions:

\begin{itemize}
    \item \textbf{RQ1:} Can we reliably estimate which inputs require more powerful models
    for AI-generated media detection, including under generator distribution shift?
    \item \textbf{RQ2:} Can such estimates be used to guide model selection across heterogeneous detectors without relying on calibrated confidence scores?
    \item \textbf{RQ3:} Does disagreement-based routing reduce verifier cost while recovering accuracy under deployment shift to unseen generator families?
\end{itemize}

\subsection{Contributions}

This work makes the following contributions:

\begin{itemize}
    \item \textbf{C1:} We show that the disagreement between artifact-based and semantic detectors identifies inputs where cheap detectors are less reliable, across generator families, and remains informative under deployment shift to unseen generators.
    
    \item \textbf{C2:} We propose a parameter-free routing strategy that uses disagreement as a difficulty signal for selective model invocation, requiring no confidence
    calibration or policy training.
    
    \item \textbf{C3:} On deployment shift benchmarks (SDXL and FLUX), cascade escalation to GPT-5.5 recovers up to 7.9 percentage points of accuracy over a stale semantic probe while reducing verifier invocations by 58--64\%. A uniform always-escalate policy, evaluated on the full benchmark, shows this closes most but not all of the achievable gap (3--7.8 percentage points remain), making the accuracy-cost tradeoff explicit rather than presenting selective escalation as a cost-free win.

    \item \textbf{C4:} We show that a learned escalation classifier (a supervised logistic-regression router over detector-derived features) converges to the disagreement rule across multiple datasets, independently validating disagreement as the dominant routing signal while providing no measurable improvement over the parameter-free rule. Against a joint artifact-and-semantic confidence-routing baseline with thresholds fit on held-out data, the disagreement rule is more accurate on two of three in-domain benchmarks but has no accuracy advantage under deployment shift, where the calibrated confidence rule is tied or better; disagreement's advantage is that it needs no threshold or calibration data.

    \item \textbf{C5:} We introduce a verifier validation gate with a prespecified, paired-bootstrap significance test that determines whether disagreement-based escalation should be enabled before deployment, reducing harmful escalations by roughly an order of magnitude relative to a point-estimate rule. Evaluated under genuine calibration/test separation, the gate reaches at least 90\% test-set agreement with the oracle decision in 21 of 32 eligible configurations at $N=100$ and 26 of 32 at $N=200$; a three-way detector selector is equivalent to the binary gate as a decision procedure.

    \item \textbf{C6:} We show that BigGAN's near-ceiling in-domain artifact-detector accuracy is substantially explained by a native-resolution shortcut rather than genuine generative-artifact detection: a classifier using only image width and height matches the trained detector's accuracy, and resolution-matching the two classes collapses the trained detector's accuracy to chance.
\end{itemize}


\section{Related Work}
This section reviews prior work related to AI-generated media detection and adaptive inference. We first identify gaps in existing approaches that motivate this work, then discuss detection methods under distribution shift, examine approaches for selective computation, and finally compare routing strategies based on confidence and learned policies.




\subsection{Gaps in Existing Approaches}

Existing forensic detection and cost-aware inference literature leaves several gaps that motivate this work.

\textbf{Single signal reliance.} Most forensic detectors rely on a single signal type, such as diffusion-based artifacts or reconstruction inconsistencies~\cite{Corvi_2023_ICASSP} or semantic embeddings from large pretrained models~\cite{cozzolino2023raising}. These approaches work well within specific generator families but degrade under generator distribution shift. Recent large-scale evaluations show that no single detector consistently generalizes across modern image generation pipelines~\cite{epstein2023online, cozzolino2023raising}.

\textbf{Confidence-based routing under shift.} Cost-aware inference methods and selective classification approaches rely on calibrated confidence estimates to decide when to defer or escalate predictions. Predictive uncertainty in deep models degrades significantly under distribution shift~\cite{ovadia2019can}. As a result, confidence-based routing signals become unreliable in the very regimes where robust detection is most needed. Existing literature on uncertainty estimation and selective prediction does not address this failure mode in heterogeneous detector cascades.

\textbf{Disagreement as a routing signal.} Inter-model disagreement has been studied extensively in active learning~\cite{seung1992query} and ensemble-based uncertainty estimation~\cite{lakshminarayanan2017simple}. In these settings, disagreement identifies informative samples or improves predictive uncertainty. Prior work treats disagreement as a mechanism for aggregation or data selection. Its role as a routing signal for cost-aware inference remains unexplored in forensic detection pipelines. Specifically, using disagreement as a diagnostic to indicate when low-cost detectors should not be trusted has not been studied.

This work addresses these gaps in four ways. First, we combine heterogeneous artifact-based and semantic signals through a parameter-free disagreement rule. Second, we demonstrate empirically that disagreement provides a more reliable routing signal than confidence under distribution shift. We evaluate this at matched operating points across multiple generator families. Third, we account for both latency and token cost when evaluating VLM-based escalation strategies. Fourth, we reinterpret disagreement as a diagnostic for selective invocation of more expensive models, not as a voting mechanism.

\textbf{Dataset shortcuts.} Published analysis has shown that detectors trained on
GenImage-family benchmarks can exploit JPEG compression and image-resolution biases
rather than genuine generative artifacts, since real and synthetic images in these
benchmarks are frequently acquired or post-processed differently. We directly measure
this concern for our own artifact detector rather than treat it as an external
literature point: Section~\ref{sec:three_cases}'s Limitations discussion reports a
native-resolution shortcut specific to the BigGAN subset (every sampled fake image is
natively $128{\times}128$ pixels, unlike the variable-resolution real class), using a
resolution-only control classifier and a resolution-matching ablation to isolate how
much of the detector's reported accuracy the shortcut explains.

\subsection{Detection under Distribution Shift}

A large body of work has focused on detecting AI-generated images using artifact-based features. Early approaches such as CNN-based detectors~\cite{wang2020cnn} exploit low-level patterns introduced by generative models. More recent methods aim to improve generalization across generators, such as universal fake detectors~\cite{ojha2023fakedetect}. Benchmark datasets such as GenImage~\cite{NEURIPS2023_f4d4a021} highlight the difficulty of this problem, showing that many detectors degrade significantly when evaluated on unseen generators.

Other approaches leverage higher-level representations. Models such as CLIP~\cite{pmlr-v139-radford21a} learn semantic features from large-scale image-text data and have been used in downstream detection tasks. Diffusion-specific detection methods, such as DIRE~\cite{Wang_2023_ICCV}, further illustrate the diversity of signals used in modern detection systems. Despite these advances, no single model consistently generalizes across generator families, motivating the use of multiple complementary detectors.

Frequency-domain analysis provides a complementary artifact signal. Spectral methods
exploit periodic and structured artifacts introduced by GAN upsampling and diffusion
denoising pipelines, and have shown strong in-domain generalization~\cite{Corvi_2023_ICASSP}.
Beyond CLIP, recent vision encoders such as DINOv2~\cite{oquab2023dinov2} and
SigLIP2~\cite{tschannen2025siglip2} offer stronger visual representations and have
demonstrated improved robustness under limited training data. Their behavior under
generator distribution shift in forensic detection tasks remains underexplored, motivating
their inclusion as semantic detector candidates in this work.


\subsection{Selective Computation and Model Cascades}

Selective computation has been studied as a way to reduce inference cost by adapting computation to input difficulty. Early work on selective classification~\cite{NIPS2017_4a8423d5} proposes abstaining on uncertain predictions to improve reliability. More recent systems extend this idea to model selection, routing inputs across models of varying capacity. For example, FrugalGPT~\cite{chen2024frugalgpt} reduces cost by selecting among large language models based on estimated utility. These approaches typically rely on confidence scores or calibrated probabilities to guide routing decisions. However, such signals can be unreliable under distribution shift, where models may produce overconfident incorrect predictions.

\subsection{Routing Signals for Model Selection}

Most existing routing strategies use confidence-based signals, such as entropy or prediction probability, to determine whether additional computation is required. While effective in controlled settings, these signals often degrade under distribution shift due to poor calibration. Learned routing policies, such as bandit-based approaches, attempt to improve decision-making by training a model to predict the best action
from input features.

In contrast, this work explores disagreement between heterogeneous detectors as a routing signal. Rather than relying on calibrated confidence scores or learned policies, we use prediction disagreement between artifact-based and semantic models to estimate input difficulty. This provides a simple, calibration-free signal that remains effective under distribution shift.

\begin{table*}[h]
\centering
\small % Reduces font size to help fit more content comfortably
\caption{Comparative analysis of related AI-generated media detection approaches. \checkmark~indicates capability is present; \texttimes~indicates absent.}
\label{tab:litsurvey}
\setlength{\tabcolsep}{4pt} % Slightly tighter spacing between columns

% Define a custom column type 'Y' that is a ragged-right version of 'X'
\newcolumntype{Y}{>{\RaggedRight\arraybackslash}X}

\begin{tabularx}{\textwidth}{p{2.0cm}YYYYYY}
\toprule
\textbf{Feature}
  & \textbf{CNNSpot} \cite{wang2020cnn}
  & \textbf{DIRE} \cite{Wang_2023_ICCV}
  & \textbf{UnivFD} \cite{ojha2023fakedetect}
  & \textbf{BusterX} \cite{wen2025busterxmllmpoweredaigeneratedvideo}
  & \textbf{FrugalGPT} \cite{chen2024frugalgpt}
  & \textbf{Ours} \\
\midrule
Detection signal
  & (\checkmark) Frequency \& upsampling artifacts
  & (\checkmark) Diffusion reconstruction error
  & (\checkmark) CLIP embeddings for generalizable detection
  & (\checkmark) MLLM open-ended visual reasoning
  & (\texttimes) LLM text output, not a forgery detector
  & (\checkmark) Artifact + semantic embeddings + VLM reasoning \\
\midrule
Model architecture
  & (\texttimes) Single CNN/fixed computation
  & (\texttimes) Single diffusion model/fixed computation
  & (\texttimes) Single CLIP backbone/fixed computation
  & (\texttimes) Single MLLM/fixed computation
  & (\checkmark) LLM cascade/variable computation
  & (\checkmark) Heterogeneous three-stage pipeline \\
\midrule
Sequential routing
  & (\texttimes) No routing/uniform inference
  & (\texttimes) No routing/uniform inference
  & (\texttimes) No routing/uniform inference
  & (\texttimes) No routing/uniform inference
  & (\checkmark) Budget-aware cascade routing
  & (\checkmark) Disagreement-based escalation (primary) \\
\midrule
Cost awareness
  & (\texttimes) Fixed cost per sample
  & (\texttimes) Fixed cost per sample
  & (\texttimes) Fixed cost per sample
  & (\texttimes) Fixed cost per sample
  & (\checkmark) Minimizes LLM API cost
  & (\checkmark) Latency and token-aware cascade \\
\midrule
Multi-model disagreement
  & (\texttimes) Single model/no signal
  & (\texttimes) Single model/no signal
  & (\texttimes) Single model/no signal
  & (\texttimes) Single model/no signal
  & (\texttimes) Not applicable to forgery detection
  & (\checkmark) Artifact/semantic detector disagreement as primary signal \\
\midrule
Generator shift eval.
  & (\texttimes) GAN only, degrades on diffusion
  & (\checkmark) Evaluated across diffusion
  & (\checkmark) Evaluated across multiple models
  & (\checkmark) Modern video generators
  & (\texttimes) Not evaluated on media forgery
  & (\checkmark) Video, GAN, diffusion, SDXL, and FLUX generators \\
\midrule
Media type
  & (\texttimes) Images only
  & (\texttimes) Images only
  & (\texttimes) Images only
  & (\checkmark) Video forgery detection
  & (\texttimes) Text only
  & (\checkmark) Video frames and static images \\
\bottomrule
\end{tabularx}
\end{table*}


\section{Proposed Solution}
This section describes the proposed framework for difficulty-aware model selection. We first formalize the routing problem and define the cost model, then introduce the disagreement-based routing strategy, examine the alternative routing baselines used for comparison, and finally detail how latency is measured across routing strategies.

\subsection{Problem Formulation}

We consider the problem of detecting AI-generated media using a set of heterogeneous detectors with different computational costs and failure modes. Given an input sample $x$, the goal is to select an appropriate subset of detectors to apply such that prediction accuracy is maintained while minimizing computational cost.

Let:
\[
\mathcal{M} = \{f_{\text{art}}, f_{\text{sem}}, f_{\text{vlm}}\}
\]
denote the set of available detectors, corresponding to an artifact-based model, a semantic model, and a high-capacity vision-language model. Each model produces a binary prediction $\hat{y}_i \in \{0,1\}$. A routing policy $\pi(x) \in \{0,1\}$ determines whether the VLM is invoked on input $x$, with the artifact and semantic detectors always applied as
low-cost first-stage filters. We evaluate multiple instantiations of this framework in which the artifact detector is instantiated as a dual-stream FFT probe (primary) or CNNSpot
(baseline), and the semantic detector as CLIP, SigLIP2, or DINOv2.

We report on two complementary cost metrics. Wall-clock latency, measured in milliseconds on a fixed GPU, applies to every detection stage:
\[
C_{\text{lat}}(x, \pi) = \kappa^{\text{lat}}_{\text{art}} + \kappa^{\text{lat}}_{\text{sem}}
+ \mathbbm{1}[\pi(x)=1] \cdot \kappa^{\text{lat}}_{\text{vlm}}(x),
\]
where $\kappa^{\text{lat}}_{\text{vlm}}(x)$ is measured per sample because VLM latency depends on input resolution. The token cost, measured as mean VLM input and output tokens per sample, applies only when the VLM is invoked and reflects the standard unit of cost for commercial VLM APIs~\cite{chen2024frugalgpt}:
\[
C_{\text{tok}}(x, \pi) = \mathbbm{1}[\pi(x)=1] \cdot \bigl(\tau^{\text{in}}_{\text{vlm}}(x)
+ \tau^{\text{out}}_{\text{vlm}}(x)\bigr).
\]
The routing objective is to minimize the expected cost under either metric while maintaining a target accuracy $\alpha$:
\[
\min_{\pi} \; \mathbb{E}_x[C(x, \pi)] \quad \text{s.t.} \quad \text{Accuracy}(\pi) \geq \alpha.
\]
We report $C_{\text{lat}}$ for every evaluated method and $C_{\text{tok}}$ wherever the verifier is an API model (GPT-5.5; Section~\ref{sec:cost_accounting}), as the deployment scenarios vary between self-hosted (latency-bound) and API-based (token-bound) inference.

\subsection{Disagreement as a Difficulty Signal}

We estimate the input difficulty using disagreement between the artifact-based and semantic detectors:

\begin{equation}
d(x) = \mathbbm{1}\left[\hat{y}_{\text{art}}(x) \neq \hat{y}_{\text{sem}}(x)\right].
\label{eq:disagree}
\end{equation}

When the detectors agree, predictions are highly reliable. Disagreement corresponds to
cases where predictions conflict, indicating that neither signal alone is sufficient and
that additional computation may be required. Unlike confidence-based measures such as
entropy, this signal requires no calibration and shows consistent diagnostic behavior
across generator families. This behavior comes from the complementary failure modes of
both artifact and semantic detectors, whose errors are weakly correlated under
distribution shift. This leads to disagreement reliably identifying the cases where cheap detectors
are not trusted.


\subsection{Disagreement-Based Routing}

Using the disagreement signal, we define a simple routing policy:

\begin{equation}
\pi_{\text{disagree}}(x) =
\begin{cases}
0, & \text{if } d(x) = 0 \\
1, & \text{if } d(x) = 1
\end{cases}
\label{eq:pidisagree}
\end{equation}

where $\pi(x) = 1$ indicates that the VLM is invoked. Inputs where the detectors agree are handled by the cheap stage alone while disagreement cases are escalated to the VLM. This enables selective allocation of computation based on estimated difficulty. Note that outputting the
semantic detector's prediction on all inputs is equivalent to this routing policy without VLM escalation, since the artifact and semantic detectors produce identical predictions on
agreement cases by definition. In-domain, this semantic-only baseline achieves the best
cost-accuracy tradeoff. The cascade's value is under deployment shift, where a stale
semantic probe loses accuracy and VLM escalation on disagreement cases recovers it.

\subsection{Generalized K-Stage Formulation and Three Cascade Cases}
\label{sec:kstage}

The proposed detector is a special case of a $K$-stage cascade. Let
$f_1, \ldots, f_K$ denote detectors with non-decreasing cost
$C_1 \le C_2 \le \cdots \le C_K$ and accuracies $A_1, \ldots, A_K$ on a
fixed evaluation distribution. The useful quantity is the conditional gap
on the inputs that stage $i$ actually routes forward. The global gap
$A_{i+1}-A_i$ is only a coarse summary. Escalation improves accuracy only
when the next stage is better on that routed subset.

In our setting, let $K=3$: $f_1 = f_{\text{art}}$, $f_2 = f_{\text{sem}}$, and
$f_3 = f_{\text{vlm}}$. The routed subset is the disagreement set
$D = \{x : \hat y_{\text{art}}(x) \neq \hat y_{\text{sem}}(x)\}$
(Eq.~\eqref{eq:disagree}). The resulting cascade accuracy is
\begin{equation}
A_{\text{cascade}} = (1-p_D)\, A_{\text{agree}} + p_D\, A_{3\mid D},
\label{eq:decomp}
\end{equation}
where $p_D=P(x\in D)$, $A_{\text{agree}}$ is accuracy on the agreement set,
and $A_{3\mid D}$ is verifier accuracy on $D$. On the agreement set,
$f_1$ and $f_2$ make the same prediction by construction
(Eq.~\eqref{eq:pidisagree}). The gain over the semantic-only detector is
therefore
\begin{equation}
\text{Gain} = A_{\text{cascade}} - A_2
= p_D \left( A_{3\mid D} - A_{2\mid D} \right),
\label{eq:gain}
\end{equation}
where $A_{2\mid D}$ is the accuracy of the semantic detector on the same routed
subset. Thus, a stronger verifier helps only if it is stronger on
disagreement cases. Average accuracy alone does not determine the cascade gain.
This leaves us with three cascade cases.

\textbf{Case 1 (both cheap stages strong, $A_1 \approx A_2 \approx A_3$).}
The artifact and semantic detectors usually agree, so $p_D$ is small. Even
a strong verifier has little effect on end-to-end accuracy because it is
rarely invoked.

\textbf{Case 2 (asymmetric cheap stages, $A_1 \ll A_2$, and $A_2 \gtrsim A_3$).}
The semantic detector is already the better cheap model. Many disagreement
cases are therefore inputs where the artifact detector is wrong and the
semantic detector is right. If $A_{3\mid D} < A_{2\mid D}$, routing replaces
correct semantic predictions with noisier verifier predictions, so
Gain $<0$. In this case, any improvement over the \emph{artifact}-only
baseline comes from the semantic detector, rather than from escalation.

\textbf{Case 3 (both cheap stages degraded, $A_1 \ll A_3$ and $A_2 \ll A_3$).}
Both cheap detectors are unreliable with disagreement being common. If the
verifier remains accurate on $D$, then $A_{3\mid D}$ can substantially
exceed $A_{2\mid D}$. This is the cascade case where Eq.~\eqref{eq:gain} becomes
large and positive, and the gain is caused by routing hard cases to the
verifier.

\textbf{Breakdown condition.} Disagreement is useful only when the two cheap
detectors fail differently and the verifier is better on the routed subset.
If $f_1$ and $f_2$ share a correlated blind spot, they can agree while both
are wrong. Those inputs never enter $D$, so the cascade cannot route them to
the verifier. In the extreme, $p_D \to 0$ and
Eq.~\eqref{eq:decomp} reduces to
$A_{\text{cascade}} \approx A_{\text{agree}}$, even when the agreement set
contains systematic errors. Section~VI maps these cases and this breakdown
condition to the empirical results in Section~V.
Section~VI.I formalizes a validation gate that estimates, from a small
labeled calibration sample rather than full post-hoc evaluation, which of
these cases a given deployment setting falls into.

\subsection{Baseline Routing Strategies}

To evaluate the effectiveness of disagreement-based routing, we compare two alternative strategies.

\textbf{Entropy-based routing.}  
We use the entropy of the artifact detector as a measure of uncertainty:

\begin{equation}
H(x) = - \sum_{i} p_i(x) \log p_i(x),
\end{equation}

where $p_i(x)$ are the predicted class probabilities. Inputs with entropy above a threshold are escalated to the VLM. This represents a standard confidence-based routing approach.

\textbf{Learned escalation classifier.}
\label{sec:bandit_rename}
An earlier version of this work described this baseline as a contextual bandit.
A reviewer correctly noted that a logistic regression trained to predict an
escalation decision, without an explicit reward signal, an exploration policy, or
an online or batch reinforcement-learning training procedure, does not establish
that a bandit was implemented. We therefore describe it plainly: a supervised
logistic-regression classifier trained on features derived from detector output,
including artifact detector confidence, semantic confidence, and prediction
disagreement, predicting whether an input should be escalated to the VLM as a
learned routing baseline. This relabeling changes no numerical result; it only
corrects the method's name.

\subsection{Latency Measurement}

Inference latency for each model component is measured independently using synchronized CUDA execution on a fixed hardware setup. The total latency per sample is the sum of the costs of all models invoked by the routing policy. Artifact and semantic detectors are modeled as additive costs, which is a  conservative estimate since these stages could run in parallel in practice. All latency measurements are intended for relative comparison across routing strategies rather than as absolute deployment guarantees.

\FloatBarrier
\begin{figure*}[h!]
    \centering
    \includegraphics[width=0.7\textwidth]{workflowfinal.png}
    \caption{Comparison of detection workflows. (a) FrugalGPT introduces cost-aware cascade
routing for homogeneous LLM outputs. (b) Our framework adapts this paradigm to
heterogeneous forensic detectors, using artifact-semantic disagreement as the primary
routing signal.}
    \label{fig:workflowdiagram}
\end{figure*}
\FloatBarrier

\section{Experiments}

This section describes the evaluation methodology used to assess the proposed routing framework. We first define the compared methods and evaluation metrics, then summarize implementation details and latency measurements.

\subsection{Evaluation Strategy}
We evaluate our framework against artifact-only and semantic-only detectors applied without routing, an always-escalate pipeline that invokes the verifier on every sample, and entropy-based and learned-classifier routing baselines. UnivFD~\cite{ojha2023fakedetect} is included
as a published 2023 baseline using CLIP ViT-L/14 features trained on GAN images. We
evaluate two verifiers: Qwen2.5-VL-7B~\cite{bai2025qwen25vltechnicalreport} as the
baseline and GPT-5.5~\cite{openai2026gpt55} as the primary, to test whether a stronger verifier better exploits the disagreement signal.
Evaluation focuses on detection accuracy, balanced accuracy, 
average computational cost, and escalation rate. For video 
datasets, a representative frame is sampled from each clip 
and treated as an image detection task following standard 
evaluation practice~\cite{wen2025busterxmllmpoweredaigeneratedvideo}.



\subsection{Implementation}

The pipeline integrates three heterogeneous detection components. The first stage is an artifact-based detector. The primary artifact detector is a dual-stream frequency-domain probe that combines radial FFT magnitude profiles~\cite{dzanic2020fourier} with DCT high-frequency statistics~\cite{frank2020leveraging, Corvi_2023_ICASSP} and color channel
coherence features, trained as a logistic regression classifier. Each image is resized
to $256{\times}256$ and converted to grayscale for the first two streams, giving a
43-dimensional feature vector: (i) a 32-bin radial profile of the log-magnitude of the
centered 2-D FFT (mean over each normalized-radius bin), L2-normalized; (ii) eight
statistics (mean, standard deviation, skewness, kurtosis, and the 5th, 25th, 75th and
95th percentiles) of the bottom-right $3{\times}3$ coefficients of every non-overlapping
$8{\times}8$ block DCT; and (iii) the three pairwise Pearson correlations between the R,
G and B channels. Features are standardized and classified with
\texttt{LogisticRegression(C=0.1, max\_iter=2000, random\_state=42)}.
CNNSpot~\cite{wang2020cnn} is retained as an artifact detector baseline. The second stage is a semantic detector. Visual embeddings are extracted with one of three backbones: CLIP
ViT-B/32~\cite{pmlr-v139-radford21a}, SigLIP2~\cite{tschannen2025siglip2}, or
DINOv2~\cite{oquab2023dinov2}. It's then passed through a logistic regression classifier trained on the first 50\% of each dataset. All three semantic probes use the same training split and evaluation protocol. The third and final stage is an optional VLM verifier which is only invoked 
on disagreement cases.


To utilize GPT-5.5 as the primary image verifier, we prompt it as follows:

\noindent
\begin{minipage}{\linewidth}
\vspace{10pt}
\itshape
``Analyze this image carefully. Determine whether it is a REAL photograph or an AI-generated FAKE image. Consider both possibilities equally. Answer with exactly one word: REAL or FAKE.''
\vspace{10pt}
\end{minipage}

When applied to the real images bundled with AIGIBench (PASCAL VOC 2007--2012
photographs), GPT-5.5 classified 82--97\% of them as \emph{FAKE} despite being
genuine photographs. PASCAL VOC images are archival scans with heavy JPEG
compression that a frontier VLM trained on modern high-quality imagery does not
recognize as real, reducing real-class recall to near zero. We replaced these
images with ImageNet val/nature photographs from the GenImage SD~1.4
partition~\cite{NEURIPS2023_f4d4a021}, which GPT-5.5 correctly identifies as real
and which match the real-image source used to train the SD~1.4 semantic probes.

Selecting a replacement after observing a model failure on the original benchmark
is a legitimate methodological concern. We therefore report the full cascade
(not only GPT-5.5's real-class recall) on three real-image sources for SDXL and
FLUX: the original bundled PASCAL VOC photographs, the ImageNet-sourced
replacement used throughout the remainder of this paper, and an independently
selected source (COCO val2017~\cite{lin2014microsoft}, Flickr-acquired, CC-licensed,
sharing no images with either PASCAL VOC or the ImageNet pool, and with
compression and resolution characteristics distinct from both). Table~\ref{tab:real_image_sources}
reports cascade accuracy under all three. Under the original PASCAL VOC images,
cascade accuracy collapses to 72.9--77.9\%, driven by an agreement-path accuracy
of only 43--53\%: the local artifact detector alone misclassifies 83--85\% of these
photographs as fake, essentially matching GPT-5.5's own 82--97\% misclassification
rate despite using an entirely different (frequency-domain, non-VLM) signal. This
shows the PASCAL VOC images degrade the \emph{entire} cascade, not only the GPT-5.5
verifier, and that the substitution was not selected to flatter a single model's
result. Under the independently chosen COCO source, cascade accuracy is
89.6--94.1\%, closely tracking the ImageNet-sourced benchmark used elsewhere in
this paper and confirming the main escalation-helps finding is not an artifact of
the specific ImageNet real-image pool. We do not recommend categorically replacing
compressed archival photographs in all VLM evaluations of forensic detectors;
genuine compressed imagery remains a relevant deployment input, and the result
above should instead be read as evidence that this particular PASCAL VOC-sourced
benchmark carries a pre-existing confound affecting any detector tested on it,
VLM or not.

\begin{table}[h]
\centering
\footnotesize
\caption{Cascade accuracy on SDXL/FLUX under three real-image sources. Agree.\ Acc.\ is agreement-path accuracy (both cheap detectors agree); n=4{,}000 per cell (2{,}000 real, 2{,}000 fake).}
\label{tab:real_image_sources}
\begin{tabular}{llcccc}
\toprule
Source & Dataset & Sem. & Agree.\ Acc. & Casc.\ Acc. & $\Delta$ vs.\ stale \\
\midrule
PASCAL VOC (original) & SDXL & SigLIP2 & 0.533 & 0.767 & +0.052 \\
PASCAL VOC (original) & SDXL & DINOv2  & 0.532 & 0.768 & +0.055 \\
PASCAL VOC (original) & FLUX & SigLIP2 & 0.525 & 0.779 & +0.025 \\
PASCAL VOC (original) & FLUX & DINOv2  & 0.431 & 0.729 & +0.052 \\
\midrule
COCO (independent)    & SDXL & SigLIP2 & 0.888 & 0.927 & +0.067 \\
COCO (independent)    & SDXL & DINOv2  & 0.886 & 0.929 & +0.090 \\
COCO (independent)    & FLUX & SigLIP2 & 0.921 & 0.941 & +0.041 \\
COCO (independent)    & FLUX & DINOv2  & 0.844 & 0.896 & +0.090 \\
\midrule
ImageNet (used throughout) & SDXL & SigLIP2 & -- & 0.925 & +0.061 \\
ImageNet (used throughout) & SDXL & DINOv2  & -- & 0.930 & +0.079 \\
ImageNet (used throughout) & FLUX & SigLIP2 & -- & 0.940 & +0.034 \\
ImageNet (used throughout) & FLUX & DINOv2  & -- & 0.893 & +0.075 \\
\bottomrule
\end{tabular}
\end{table}


Confidence scores are derived from the logistic probe output for the FFT artifact detector and all semantic detectors. For the CNNSpot baseline, confidence is the maximum softmax probability from the ResNet output. The learned escalation classifier and entropy baselines use these scores as features. The disagreement rule uses only binary predictions and requires no confidence
calibration.

The routing policy operates on the joint output of the artifact and semantic detectors. 
In the disagreement threshold configuration, escalation is triggered whenever the two first-stage detectors produce conflicting 
predictions. In the entropy threshold configuration, escalation is
triggered when the artifact detector's prediction entropy exceeds $\tau = 0.5$. The
learned escalation classifier (a logistic-regression router, not a reinforcement-learning
bandit -- Section~\ref{sec:bandit_rename}) is trained on artifact
confidence, semantic confidence, and their disagreement signal to produce
a binary escalation decision.

We measure deployment latency on an NVIDIA RTX A4500 GPU with synchronized CUDA execution. The FFT probe takes 2.5~ms per sample. Each semantic detector takes 10.4~ms per sample,
giving a cheap-stage latency of 12.9~ms. This represents a 4.4$\times$ reduction in cheap-stage cost compared to the CNNSpot-based configuration (56.2~ms). GPT-5.5 latency, measured as mean per-call wall-clock time over the current, fully-corrected disagreement-subset predictions (Section~\ref{sec:gpt55_audit}), averages 4{,}526~ms on GenBuster, 4{,}691~ms on SD~1.4, and 3{,}827~ms on BigGAN. On the SDXL and FLUX benchmarks, GPT-5.5 latency averages 3{,}040--3{,}102~ms and 4{,}509--4{,}778~ms respectively. These figures are markedly higher than an earlier draft's reported 1{,}595/1{,}459/1{,}572~ms (in-domain) and 3{,}670--3{,}849~ms (FLUX): the earlier numbers were measured against GPT-5.5 calls made before the output-token-budget fix described in Section~\ref{sec:gpt55_audit}, which capped the model's response at 16--100 tokens and frequently truncated its reasoning; the corrected calls, which escalate the budget to let the model complete a reasoning pass before answering, take substantially longer per call, and that cost is reported here rather than the faster but incorrect pre-fix figures. Qwen2.5-VL
latency on the in-domain benchmarks is 190~ms (GenBuster), 83~ms (SD~1.4), and 71~ms
(BigGAN) under always-escalate. We report both latency and token cost throughout our results (Section~\ref{sec:cost_accounting}). Latency comparisons assume sequential FFT and semantic detector execution. Parallel
execution would further reduce cheap-stage cost.

\subsubsection{GPT-5.5 and Qwen Verifier Audit}
\label{sec:gpt55_audit}

Every GPT-5.5 call in this paper uses the \texttt{gpt-5.5} model identifier via the
OpenAI Responses API, with no explicit \texttt{reasoning\_effort}, \texttt{verbosity},
or \texttt{temperature} parameter set; these are left at the API's own defaults, which
we disclose rather than imply a specific configured value. An earlier version of the
calling code capped the output-token budget at 16--100 tokens; because \texttt{gpt-5.5}
is a reasoning model that spends output-token budget on a hidden reasoning pass before
emitting a visible answer, this cap frequently produced empty or mid-word-truncated
responses (e.g.\ \texttt{"FA"} instead of \texttt{"FAKE"}), which earlier code silently
defaulted to a fixed label rather than treating as invalid. The corrected implementation
escalates the output-token budget across up to 5 retries ($2{,}048 \to 4{,}096$ tokens,
exponential backoff), parses a response as valid only if it is exactly \texttt{"REAL"}
or \texttt{"FAKE"} after case-normalization, and never guesses a label for an invalid or
empty response; such responses are recorded as missing rather than defaulted. We
identified and corrected 19 responses across the SDXL/FLUX disagreement predictions that
had been silently mislabeled by the pre-fix truncation-detection logic (non-empty but
incomplete text, e.g.\ \texttt{"FA"} at exactly the old 100-token cap, which earlier code
did not recognize as invalid), together with the disagreement-subset backfill rows used
to complete Table~\ref{tab:verifier_comparison}. GPT-5.5 coverage of the disagreement
subset is complete (100\%) for the uniform-escalation and real-image-source results
reported in Sections~\ref{sec:real_image_dedup} and~\ref{sec:cost_accounting}, which were
computed via the OpenAI Batch API (18{,}501 total calls across both verifiers' corrected
code path, 18{,}490 valid on first pass, the remainder recovered via a small synchronous
follow-up); it is partial (93--95\%) for the Qwen columns of Table~\ref{tab:e8_results},
where uncovered disagreement samples retain the stale semantic probe's prediction rather
than being dropped, and partial for the GenBuster/SD~1.4 cells of
Table~\ref{tab:verifier_comparison} as already noted there. We do not have per-call
timestamps for GPT-5.5/Qwen predictions made before this correction (file modification
times are the only available proxy, clustering within a roughly 9-hour window); all
calls made as part of this correction carry explicit UTC ISO-8601 timestamps, retry
counts, and full raw response text, released alongside the per-sample prediction files.

\section{Results}

This section presents empirical results across video, diffusion, and GAN benchmarks. We first describe the datasets and evaluation protocol, then compare detector performance, analyze routing behavior, and finally study the robustness of the proposed signal through ablations.

\subsection{Experimental Setup}
\label{sec:exp_setup}

We evaluate on three in-domain benchmarks spanning distinct generator families.
GenBuster-200K-mini~\cite{wen2025busterxmllmpoweredaigeneratedvideo} comprises 1{,}000
samples from four modern AI video generators and 1{,}000 real video clips.
GenImage SD~1.4~\cite{NEURIPS2023_f4d4a021} comprises 6{,}000 Stable Diffusion v1.4
images and 6{,}000 real ImageNet images. GenImage
BigGAN~\cite{NEURIPS2023_f4d4a021} comprises 6{,}000 BigGAN images and 6{,}000 real
ImageNet images. Each sample is represented by a single frame or image following standard
practice~\cite{wen2025busterxmllmpoweredaigeneratedvideo}.

We additionally evaluate under deployment shift on two modern diffusion benchmarks.
Fake images are drawn from AIGIBench~\cite{li2025artificial} SDXL and FLUX1-dev
subsets. Real images are ImageNet val/nature photographs from the GenImage SD~1.4
partition. AIGIBench's bundled PASCAL VOC real images are not used because they
cause systematic VLM calibration failure (Section~IV.B). Each benchmark contains
4{,}000 samples (2{,}000 fake, 2{,}000 real). All semantic probes are trained
exclusively on SD~1.4 data. No SDXL or FLUX images appear in any training split.


For the in-domain benchmarks, each dataset is partitioned into exactly two
non-overlapping, stratified splits: a fixed-size evaluation set, and a training
set consisting of every remaining sample in the pool. Evaluation sets are drawn by
stratified (real/fake-balanced) sampling at a fixed target size per dataset: 1{,}000
samples for GenBuster (the dataset's full pool is only 2{,}000 samples, so this is a
disclosed ceiling rather than a fraction of a larger pool) and 2{,}000 samples each for
SD~1.4 and BigGAN (each drawn from a 12{,}000-sample pool). All detectors --
the FFT artifact probe, the semantic probes (CLIP, SigLIP2, DINOv2), and CNNSpot --
are trained on the complementary training split (the full pool minus the evaluation
set: 1{,}000 samples for GenBuster, 10{,}000 for SD~1.4 and BigGAN) using the
identical split for every detector, so all reported comparisons are evaluated on the
same held-out samples. No additional nested validation holdout is used; with a fixed
logistic-regression classifier ($C=0.1$) requiring no early-stopping criterion, a
further split was unnecessary and is not part of the current pipeline. Split
identifiers (per-dataset lists of evaluation-set paths) are released alongside the
code. We also verified, by exact-path overlap checking across all five benchmarks,
that no image is shared between any dataset's own training and evaluation splits;
We separately checked for overlap between the \emph{real}-class images reused
across different benchmarks, since SDXL and FLUX's real class is drawn from the
same ImageNet pool used for SD~1.4 (Section~\ref{sec:exp_setup}).
\label{sec:real_image_dedup}
This check found that 100\% of SDXL's and FLUX's 2{,}000 real images are drawn from
the same 6{,}000-image SD~1.4 real pool, and that 31.5\% (SDXL) and 33.1\% (FLUX) of
SD~1.4's own 1{,}000-sample evaluation-set real images are also present, unmodified,
in the SDXL/FLUX real class -- i.e.\ roughly a third of the real images used to
measure in-domain SD~1.4 accuracy are reused as the real class of the
deployment-shift benchmarks reported in Section~\ref{sec:cross_family_deployment}
and Table~\ref{tab:e8_results}. This overlap does not affect any individual
detector's own training/evaluation separation (SD~1.4 and SDXL/FLUX remain
distinct evaluation sets with distinct fake-class images and distinct accuracy
numbers), but we disclose it because it means the two benchmarks' real-class
difficulty is correlated rather than fully independent. All cross-dataset
real-image overlap counts, and confirmation of zero unintended overlap between
GenBuster/BigGAN and any other benchmark, are released alongside the split
identifiers above.

We compare FFT-only, CNNSpot-only~\cite{wang2020cnn}, UnivFD~\cite{ojha2023fakedetect},
and semantic-only configurations using CLIP~\cite{pmlr-v139-radford21a},
SigLIP2~\cite{tschannen2025siglip2}, and DINOv2~\cite{oquab2023dinov2}. Routing
baselines include entropy threshold ($\tau = 0.5$), disagreement-based escalation to Qwen2.5-VL-7B or GPT-5.5, and uniform escalation to Qwen2.5-VL. All configurations
are evaluated under identical conditions on the same held-out splits.

Accuracy figures report point estimates unless confidence intervals are shown. Balanced accuracy is the mean of fake and real recall. Latency
is measured in milliseconds on an NVIDIA RTX A4500 GPU.

\subsection{Detector Performance}


\textbf{Artifact detectors.} Table~\ref{tab:main_results} compares the dual-stream
FFT probe against CNNSpot across all three in-domain benchmarks. FFT is more accurate
and substantially faster on every dataset, with its largest accuracy advantage on
BigGAN.

\textbf{Semantic detectors.} Among the three semantic probes, SigLIP2 achieves the
strongest in-domain accuracy on GenBuster and is effectively tied with CLIP on BigGAN.
CLIP is marginally ahead on SD~1.4, with substantially overlapping confidence intervals.
We use SigLIP2 as the primary semantic backbone for consistency across evaluated
conditions. All three probes have the same measured latency.

In-domain semantic-only inference provides the strongest low-cost operating point.
Qwen escalation remains below the leading semantic probes because Qwen is weaker on
the routed disagreement subsets. This is \textbf{Case 2} in the K-stage formulation
(Section~III.D). On GenBuster in particular, the routed inputs are not intrinsically
difficult for every cheap detector: the semantic probe remains 96.1\% accurate on this
same disagreement subset (Table~\ref{tab:verifier_comparison}). Disagreement here
primarily flags inputs where the artifact detector degrades while the semantic
detector remains reliable, and escalating to a weaker verifier (Qwen) replaces a
correct semantic prediction with a less reliable one rather than resolving a case
that was difficult for both cheap detectors.

Although the disagreement cascade outperforms the FFT baseline on GenBuster and
SD~1.4 (Table~\ref{tab:main_results}), this gain comes from replacing the
artifact detector with the stronger semantic detector on disagreement cases.
It should not be interpreted as evidence that VLM escalation is helping in this
setting. Escalation becomes useful only in Case 3, where the verifier is stronger
than the cheap detectors on the routed subset.

\textbf{Verifier comparison.} Table~\ref{tab:verifier_comparison} compares Qwen2.5-VL
and GPT-5.5 on the FFT+SigLIP2 disagreement subset used by the primary cascade.
GPT-5.5 coverage on GenBuster and SD~1.4 is
partial, so we do not use those cells to support a
GPT-5.5-specific conclusion. Across the complete Qwen evaluations, verifier accuracy
remains below the best cheap detector on all three benchmarks, reinforcing that
Case~2/Case~3 identity depends on the verifier, detector pair, and deployment dataset.

\textbf{UnivFD baseline.} UnivFD~\cite{ojha2023fakedetect}, trained on GAN images using
CLIP ViT-L/14, fails on GenBuster and SD~1.4 but remains effective on BigGAN. This
contrast shows that backbone strength alone does not overcome a mismatch between the
detector's training artifacts and the target generator family.



\begin{figure*}[t]
    \centering
    \includegraphics[width=.7\textwidth]{plot_accuracy_comparison.png}
   \caption{Detector accuracy across the three in-domain benchmarks (values in Table~\ref{tab:main_results}). Bars are colored by method family: artifact detectors, semantic probes, FFT+semantic disagreement cascades with Qwen2.5-VL-7B, and reference policies (artifact-entropy routing and always escalating).}

    \label{fig:accuracy_comparison}
\end{figure*}


\begin{table}[h]
\centering
\footnotesize
\setlength{\tabcolsep}{2.5pt}
\caption{Detection accuracy and latency on the in-domain evaluation sets (GenBuster $N{=}1{,}000$; SD~1.4 and BigGAN $N{=}2{,}000$ each).}
\label{tab:main_results}
\begin{tabular}{lcccccc}
\toprule
 & \multicolumn{2}{c}{GenBuster} & \multicolumn{2}{c}{SD~1.4} & \multicolumn{2}{c}{BigGAN} \\
\cmidrule(lr){2-3} \cmidrule(lr){4-5} \cmidrule(lr){6-7}
Method & Acc. & Lat. & Acc. & Lat. & Acc. & Lat. \\
\midrule
FFT          & 0.727 &  2.5 & 0.775 &  2.5 & 0.997 &  2.5 \\
CNNSpot      & 0.534 & 45.8 & 0.487 & 45.8 & 0.705 & 45.8 \\
UnivFD       & 0.518 & 56.2 & 0.518 & 56.2 & 0.910 & 56.2 \\
CLIP         & 0.943 & 10.4 & 0.933 & 10.4 & 0.992 & 10.4 \\
SigLIP2      & 0.966 & 10.4 & 0.926 & 10.4 & 0.992 & 10.4 \\
DINOv2       & 0.953 & 10.4 & 0.877 & 10.4 & 0.984 & 10.4 \\
Entropy      & 0.742 & 17.9 & 0.695 & 14.3 & 0.997 & 12.9 \\
Casc. (C)    & 0.891 & 62.0 & 0.887 & 34.2 & 0.995 & 13.4 \\
Casc. (S)    & 0.899 & 62.2 & 0.874 & 35.0 & 0.995 & 13.4 \\
Casc. (D)    & 0.895 & 64.3 & 0.851 & 37.0 & 0.991 & 14.0 \\
Always esc.  & 0.713 & 190.2 & 0.576 & 83.2 & 0.690 & 70.6 \\
\bottomrule
\end{tabular}
\end{table}
\begin{table}[h]
\centering
\footnotesize
\setlength{\tabcolsep}{2pt}
\caption{Verifier accuracy on the FFT+SigLIP2 disagreement subset.}
\label{tab:verifier_comparison}
\begin{tabular}{llcccc}
\toprule
Dataset & Verifier & Acc. & $n$ & Best-cheap & Clears? \\
\midrule
GenBuster & Qwen2.5-VL-7B  & 0.703 & 259 & 0.961 & No \\
GenBuster & GPT-5.5        & 0.865 &  52 & 0.961 & No \\
\midrule
SD~1.4    & Qwen2.5-VL-7B  & 0.586 & 532 & 0.784 & No \\
SD~1.4    & GPT-5.5        & 0.739 &  23 & 0.826 & No \\
\midrule
BigGAN    & Qwen2.5-VL-7B  & 0.571 &  14 & 0.857 & No \\
BigGAN    & GPT-5.5        & 0.500 &  14 & 0.857 & No \\
\bottomrule
\end{tabular}
\end{table}

\subsubsection{Baseline Comparison on the Disagreement Subset}
\label{sec:baseline_comparison}

To rigorously test whether the verifier is worth invoking, we compare its
accuracy on each disagreement subset against three baselines: a random
predictor (mean and 95\% CI over 10{,}000 Monte Carlo draws), a
majority-class predictor, and the best-performing cheap detector on the
same subset ($\max(A_{1\mid D}, A_{2\mid D})$, the quantity
Eq.~\eqref{eq:gain} requires the verifier to exceed). Table~\ref{tab:baselines}
summarizes results across all 35 tested verifier-on-disagreement
configurations, with significance from a paired bootstrap (10{,}000
resamples). A verifier ``clears'' a baseline when the paired-bootstrap 95\%
confidence interval on (verifier accuracy $-$ baseline accuracy) lies
entirely above zero.

\begin{table}[h]
\centering
\small
\caption{Verifier accuracy vs.\ random, majority-class, and best-cheap-detector baselines, aggregated across 35 configurations spanning five experimental settings.}
\label{tab:baselines}
\begin{tabular}{lc}
\toprule
Baseline & Configurations cleared \\
\midrule
Random          & 28 / 35 \\
Majority-class  & 22 / 35 \\
Best-cheap-detector & 6 / 35 \\
\bottomrule
\end{tabular}
\end{table}

The six configurations that clear the best-cheap-detector baseline are the
four SDXL/FLUX GPT-5.5 results below, one cross-family pairing
(BigGAN$\to$GenBuster), and one low-training-data configuration
(SD~1.4/SigLIP2 at 1\% training fraction). The remaining 29 configurations
do not, confirming that disagreement alone does not guarantee a verifier is
worth invoking. For Qwen on the SD~1.4 disagreement subset, accuracy is
58.6\% at $N{=}532$,
above the 50.0\% random baseline (95\% CI $[45.9\%, 54.3\%]$) and the
52.8\% majority-class baseline, but well short of the 78.4\% best-cheap
baseline (paired bootstrap 95\% CI on the gap: $[-0.250,-0.145]$). Qwen is
therefore better than guessing on this subset, just not better than the
detector already available at no additional cost, which is the comparison
that determines whether escalation is worthwhile.

\begin{table*}[t]
\centering
\footnotesize
\setlength{\tabcolsep}{3pt}
\caption{Deployment-shift results on SDXL and FLUX using SigLIP2 and DINOv2. Acc/Gain reports cascade accuracy and its change from stale-semantic accuracy. GPT-5.5 columns use complete (100\%) verifier coverage of the disagreement subset. Qwen columns reflect 93--95\% coverage; uncovered disagreement samples fall back to the stale semantic prediction (see Section~\ref{sec:gpt55_audit} for coverage accounting).}
\label{tab:e8_results}
\begin{tabular}{lcccc}
\toprule
Sem. & \multicolumn{2}{c}{SDXL} & \multicolumn{2}{c}{FLUX} \\
\cmidrule(lr){2-3} \cmidrule(lr){4-5}
 & Qwen (Acc/Gain) & GPT-5.5 (Acc/Gain) & Qwen (Acc/Gain) & GPT-5.5 (Acc/Gain) \\
\midrule
SigLIP2 & 0.824 / $-$0.040 & 0.925 / $+$0.061 & 0.777 / $-$0.129 & 0.940 / $+$0.034 \\
DINOv2  & 0.819 / $-$0.032 & 0.930 / $+$0.079 & 0.738 / $-$0.080 & 0.893 / $+$0.075 \\
\bottomrule
\end{tabular}
\end{table*}

Table~\ref{tab:positive_cells} audits the four GPT-5.5 cells behind this positive
result on their full disagreement subsets, all of which GPT-5.5 labeled. About two
thirds of each subset is fake. GPT-5.5 exceeds the best cheap detector on the subset by
8.1 to 20.4 points, with every paired-bootstrap interval well above zero. The smallest
margin is FLUX/SigLIP2, where the stale probe remains comparatively strong.

\begin{table}[h]
\centering
\footnotesize
\setlength{\tabcolsep}{3pt}
\caption{Audit of the SDXL/FLUX GPT-5.5 cells on the full FFT+semantic disagreement subset (100\% GPT-5.5 coverage). Best cheap is the stale semantic probe in every cell. Gap = GPT-5.5 $-$ best cheap, with a paired-bootstrap 95\% CI (10{,}000 resamples).}
\label{tab:positive_cells}
\resizebox{\columnwidth}{!}{%
\begin{tabular}{lccccc}
\toprule
Config & $n$ & Real/Fake & GPT-5.5 & Best cheap & Gap [95\% CI] \\
\midrule
SDXL/SigLIP2 & 1{,}452 & 458/994 & 0.996 & 0.828 & $+$0.168 [0.148, 0.187] \\
SDXL/DINOv2 & 1{,}539 & 507/1032 & 0.997 & 0.793 & $+$0.204 [0.184, 0.225] \\
FLUX/SigLIP2 & 1{,}675 & 490/1185 & 0.971 & 0.890 & $+$0.081 [0.065, 0.098] \\
FLUX/DINOv2 & 1{,}659 & 522/1137 & 0.967 & 0.787 & $+$0.180 [0.159, 0.202] \\
\bottomrule
\end{tabular}
}
\end{table}


\subsection{Routing Analysis}

In this section we analyze whether artifact-semantic disagreement between the FFT artifact detector and
the semantic probe provides a reliable routing signal, and how selective escalation
affects the cost-accuracy tradeoff across in-domain and deployment-shift conditions.

Table~\ref{tab:cheap_path} reports the in-domain disagreement rates. Disagreement is
limited on GenBuster and SD~1.4 and nearly absent on BigGAN, where FFT is already near
ceiling. It rises under the SDXL and FLUX shifts, producing a larger candidate subset
for escalation without guaranteeing that the verifier will improve it.

When the FFT artifact detector and semantic probe agree, the agreement subset achieves
above 96\% accuracy across all three in-domain benchmarks (Table~\ref{tab:cheap_path}).
When they disagree, FFT accuracy on the disagreement subset drops to
3.9\% on GenBuster and 21.6\% on SD~1.4. On GenBuster, SigLIP2 remains
highly accurate on this same subset (96.1\%, Table~\ref{tab:verifier_comparison}),
so disagreement there specifically flags artifact-detector failure, not
a case where both detectors fail. On SD~1.4 the two detectors degrade
together: SigLIP2's own disagreement-subset accuracy is 78.4\%, well
below its 92.6\% overall accuracy. This asymmetry
confirms that artifact-semantic disagreement is a stable routing signal without requiring
confidence calibration.

Table~\ref{tab:escalation} breaks down escalation behavior by case type on GenBuster for
the FFT+SigLIP2 configuration. Nearly every FFT error produces disagreement and is
eligible for escalation. The remaining shared errors expose the rule's fundamental blind
spot: correlated failures cannot be detected through disagreement. Their frequency also
shows why sufficiently large evaluation sets are necessary to estimate this failure mode.

\begin{table}[h]
\centering
\caption{Escalation behavior by case type on GenBuster-200K-mini using FFT+SigLIP2 ($N{=}1{,}000$).}
\label{tab:escalation}
\begin{tabular}{lcc}
\toprule
Case & $n$ & Disagree esc. \\
\midrule
FFT correct, models agree    & 717 & 0.00 \\
FFT correct, models disagree &  10 & 1.00 \\
FFT wrong, models agree      &  24 & 0.00 \\
FFT wrong, models disagree   & 249 & 1.00 \\
\bottomrule
\end{tabular}
\end{table}

Table~\ref{tab:per_generator} breaks down accuracy by generator on GenBuster.
CNNSpot is 4--17\% accurate on every fake generator, confirming poor transfer from GAN
artifacts to modern video content. FFT ranges from 61\% (LTX-Video) to 80\%
(HunyuanVideo). The cascade improves on FFT for four of the five groups and ties it on
HunyuanVideo. SigLIP2 alone is more accurate than the cascade on every generator and on
real videos, consistent with Case~2: Qwen is weaker than the semantic probe on the
routed subset.


\begin{table}[h]
\centering
\caption{Per-generator detection accuracy on the GenBuster evaluation set of Table~\ref{tab:main_results} ($N{=}1{,}000$). Cascade is FFT+SigLIP2 disagreement routing to Qwen2.5-VL-7B.}
\label{tab:per_generator}
\resizebox{\columnwidth}{!}{%
\begin{tabular}{lrcccc}
\toprule
Generator & $n$ & CNNSpot & FFT & SigLIP2 & Cascade \\
\midrule
CogVideoX     & 155  & 0.155 & 0.665 & 0.987 & 0.845 \\
EasyAnimate   & 200  & 0.060 & 0.775 & 0.985 & 0.955 \\
HunyuanVideo  & 50   & 0.040 & 0.800 & 0.860 & 0.800 \\
LTX-Video     & 95   & 0.168 & 0.611 & 0.958 & 0.842 \\
\midrule
Real          & 500  & 0.960 & 0.742 & 0.964 & 0.914 \\
\bottomrule
\end{tabular}
}
\end{table}


The learned-escalation-classifier comparison is reported in Section~V.D.4 and interpreted in
Section~VI.C.

Selective routing reduces latency substantially relative to uniform Qwen invocation
(Table~\ref{tab:main_results}). In-domain, the semantic-only baseline remains faster and
more accurate than escalation, separating routing efficiency from verifier utility.

\textbf{Joint confidence-routing baseline.} The entropy baseline in
Table~\ref{tab:main_results} thresholds only the artifact detector's confidence. As a
stronger comparison, we define a joint signal
$H_{\text{joint}}(x) = \max\big(H(p_a(x)),\, H(p_s(x))\big)$, the binary entropy of the
less confident of the two cheap probes, and escalate when $H_{\text{joint}}(x) > \tau$.
The threshold $\tau$ is the $(1-\hat{p}_D)$-quantile of $H_{\text{joint}}$ on a
validation split, where $\hat{p}_D$ is the disagreement rate on that same split, so both
rules target the same escalation rate. In-domain, the validation split is a seeded 50\%
sub-split of each dataset's probe-training pool, disjoint from the evaluation set. SDXL
and FLUX have no training split, so $\tau$ is fit on SD~1.4's validation split with the
same stale probes and then frozen. Inputs that are not escalated keep the semantic
probe's prediction. Escalated inputs receive the same verifier labels the disagreement
rule uses: Qwen2.5-VL-7B in-domain, and on SDXL/FLUX the uniform GPT-5.5 run of
Section~\ref{sec:cost_accounting}, which labels every image.

\begin{table}[h]
\centering
\footnotesize
\setlength{\tabcolsep}{2pt}
\caption{Disagreement routing versus joint confidence routing. Accuracy with escalation rate in parentheses; $\Delta$ = joint confidence $-$ disagreement, with a paired-bootstrap 95\% CI (10{,}000 resamples). $^\ast$ CI excludes zero.}
\label{tab:confidence_routing}
\resizebox{\columnwidth}{!}{%
\begin{tabular}{llcccc}
\toprule
Config & Verifier & $n$ & Disagreement & Joint conf. & $\Delta$ [95\% CI] \\
\midrule
GenBuster & Qwen-7B & 1{,}000 & 0.899 (0.259) & 0.894 (0.267) & $-$0.005 [$-$0.025, $+$0.014] \\
SD~1.4 & Qwen-7B & 2{,}000 & 0.874 (0.266) & 0.852 (0.250) & $-$0.022$^\ast$ [$-$0.037, $-$0.007] \\
BigGAN & Qwen-7B & 2{,}000 & 0.995 (0.007) & 0.988 (0.016) & $-$0.007$^\ast$ [$-$0.011, $-$0.004] \\
\midrule
SDXL/SigLIP2 & GPT-5.5 & 4{,}000 & 0.925 (0.363) & 0.938 (0.357) & $+$0.014$^\ast$ [$+$0.006, $+$0.021] \\
SDXL/DINOv2 & GPT-5.5 & 4{,}000 & 0.930 (0.385) & 0.923 (0.354) & $-$0.006 [$-$0.015, $+$0.002] \\
FLUX/SigLIP2 & GPT-5.5 & 4{,}000 & 0.940 (0.419) & 0.942 (0.318) & $+$0.001 [$-$0.005, $+$0.008] \\
FLUX/DINOv2 & GPT-5.5 & 4{,}000 & 0.893 (0.415) & 0.894 (0.335) & $+$0.002 [$-$0.008, $+$0.011] \\
\bottomrule
\end{tabular}
}
\end{table}

Table~\ref{tab:confidence_routing} shows that the disagreement rule is significantly more
accurate on SD~1.4 and BigGAN and statistically tied on GenBuster. Under the SDXL/FLUX
shift it has no accuracy advantage. The joint rule is significantly better on
SDXL/SigLIP2 and tied on the other three configurations, and on FLUX it matches the
disagreement rule's accuracy while escalating about 10 percentage points fewer inputs.
Because the frozen threshold does not transfer exactly to the evaluation set, we also
matched escalation counts exactly by escalating the $k$ highest-$H_{\text{joint}}$
inputs, where $k$ is the disagreement rule's count (no labels used). This gives the same
picture: the joint rule is significantly better in three of the four shifted
configurations (by 0.9 to 1.7 points) and tied in the fourth, while the disagreement rule
stays ahead on SD~1.4 and BigGAN. If non-escalated inputs instead keep the more confident
probe's prediction, the joint rule loses accuracy under shift, so its strength there comes
mainly from falling back to the semantic probe. Disagreement routing is therefore not a
more accurate signal than calibrated confidence routing in general. Its advantages are
practical. It needs no threshold, no calibration data, and no probability outputs, and in
no configuration is it more than 1.7 points less accurate than the calibrated joint rule.


\begin{figure*}[h]
    \centering
    \includegraphics[width=.85\textwidth]{plot_cost_accuracy_all.png}
    \caption{Effect of disagreement escalation for each semantic backbone. Left: in-domain, the cascade (filled) adds latency and lowers accuracy relative to the semantic probe alone (hollow). Right: under the SDXL/FLUX shift, escalating 36--42\% of inputs to GPT-5.5 raises accuracy over the stale probe.}
    \label{fig:cost_accuracy}
\end{figure*}


\subsection{Ablation Studies}
The ablation studies characterize cheap-path coverage, robustness of the
disagreement-based routing signal, and cascade behavior under cross-family probe
transfer.


\subsubsection{Cheap-Path Coverage}

Table~\ref{tab:cheap_path} quantifies cheap-path coverage for the FFT+SigLIP2
configuration. Most samples remain on the high-accuracy agreement path, while BigGAN
requires almost no escalation. This supports disagreement as a sparse trigger rather
than a replacement classifier.

\begin{table}[h]
\centering
\caption{Agreement-subset coverage and accuracy for FFT+SigLIP2 on the in-domain evaluation sets.}

\label{tab:cheap_path}
\begin{tabular}{lcccc}
\toprule
Dataset & Agree & Cheap Acc & Disagree & Esc. \\
        & $n$   &           & $n$      & rate \\
\midrule
GenBuster & 741  & 0.968 & 259 & 0.259 \\
SD~1.4    & 1468 & 0.978 & 532 & 0.266 \\
BigGAN    & 1986 & 0.998 &  14 & 0.007 \\
\bottomrule
\end{tabular}
\end{table}

\subsubsection{Cheap-Trigger Ablation}
\label{sec:cheap_trigger}

Since the FFT detector contributes no independent prediction on agreement cases by
construction (Section~VI.D), we test whether its value comes from the specific
disagreement signal or from escalating a matched-rate subset of inputs. We compare the
disagreement-routed cascade with a random-trigger baseline. For each dataset, the random
baseline escalates the same fraction of inputs as the FFT+SigLIP2 disagreement rule, but
chooses those inputs uniformly at random. We repeat this process over 2{,}000 draws and
report a 95\% confidence interval for the random baseline's cascade accuracy.

Table~\ref{tab:cheap_trigger} shows the gain over random triggering is
dataset-dependent rather than uniform. On SD~1.4 the gain is practically
substantial (0.874 vs.\ 0.833). On BigGAN the gain is statistically outside the
random baseline's confidence interval but small in absolute terms (0.995 vs.\
0.990), consistent with a ceiling effect since both are already near-perfect. On
GenBuster, the disagreement-routed cascade does \emph{not} outperform the random
baseline (0.899 vs.\ 0.900, within the random baseline's confidence interval) because
its semantic probe remains reliable on disagreement cases, where escalation is
already least useful (Case 2, Section~V.B); disagreement routing's value here is
selecting a more informative subset than random \emph{only} when that subset
contains cases both cheap detectors struggle with, not as a universal property of
the rule.

\begin{table}[h]
\centering
\caption{Disagreement routing versus a matched-rate random trigger on the evaluation sets of Table~\ref{tab:main_results} (2{,}000 draws; 95\% CI). Inputs the random trigger does not escalate keep the SigLIP2 prediction.}
\label{tab:cheap_trigger}
\footnotesize
\setlength{\tabcolsep}{3pt}
\resizebox{\columnwidth}{!}{%
\begin{tabular}{l c c c c}
\toprule
Dataset & Esc. Rate & Disag. Casc. & Random Casc. [95\% CI] & Beats Random? \\
\midrule
GenBuster & 0.259 & 0.899 & 0.900 [0.887, 0.914] & No \\
SD~1.4 & 0.266 & 0.874 & 0.833 [0.822, 0.844] & Yes \\
BigGAN & 0.007 & 0.995 & 0.990 [0.988, 0.992] & Yes \\
\bottomrule
\end{tabular}
}
\end{table}

\subsubsection{Probe-Degradation Sweep}
\label{sec:limited_data}

We next test how cascade gain changes as the semantic probe becomes stronger. We hold
the FFT artifact detector fixed and vary the fraction of labeled training data available
to the semantic probe (1\%, 5\%, 10\%, 25\%, 50\%, 100\%). At each fraction, we compare
semantic-only accuracy with FFT+semantic disagreement routing using Qwen as the verifier.
Table~\ref{tab:limited_data} reports the results of SD~1.4 with SigLIP2 and DINOv2
backbones.

\begin{table}[h]
\centering
\caption{Probe-degradation sweep on SD~1.4. Cascade gain is reported relative to the semantic-only probe at the same training fraction.}
\label{tab:limited_data}
\small
\begin{tabular}{lcccc}
\toprule
Backbone & Train frac. & Sem.-only & Cascade & Gain \\
\midrule
SigLIP2  & 1\%   & 0.744 & 0.804 & $+$0.060 \\
SigLIP2  & 5\%   & 0.834 & 0.839 & $+$0.005 \\
SigLIP2  & 10\%  & 0.881 & 0.861 & $-$0.020 \\
SigLIP2  & 25\%  & 0.891 & 0.860 & $-$0.031 \\
SigLIP2  & 50\%  & 0.924 & 0.875 & $-$0.049 \\
SigLIP2  & 100\% & 0.926 & 0.874 & $-$0.052 \\
\midrule
DINOv2   & 1\%   & 0.675 & 0.757 & $+$0.083 \\
DINOv2   & 5\%   & 0.759 & 0.792 & $+$0.034 \\
DINOv2   & 10\%  & 0.806 & 0.818 & $+$0.011 \\
DINOv2   & 25\%  & 0.838 & 0.832 & $-$0.006 \\
DINOv2   & 50\%  & 0.861 & 0.841 & $-$0.020 \\
DINOv2   & 100\% & 0.877 & 0.851 & $-$0.025 \\
\bottomrule
\end{tabular}
\end{table}

Both backbones move from positive to negative cascade gain as the semantic probe
improves. This matches Eq.~\eqref{eq:gain}: escalation helps only while Qwen is stronger
than the probe on the disagreement subset.

GenBuster does not cross over because its semantic probe remains stronger at the smallest
training fraction. BigGAN remains in Case~1 because its artifact detector is already near
ceiling. The transition therefore depends on relative routed-subset strength, rather than
training-set size alone.

\subsubsection{Cross-Dataset Learned-Classifier Transfer}

Table~\ref{tab:transfer} reports cross-dataset learned-escalation-classifier transfer
accuracy using the same FFT+SigLIP2+Qwen features as the disagreement rule. The
in-domain classifier closely tracks the parameter-free rule, while cross-dataset
transfer never improves on it. The learned
policy recovers the same dominant routing signal without providing a robust transfer
advantage.



\begin{table}[h]
\centering
\caption{Cross-dataset learned-escalation-classifier transfer on the evaluation sets of Table~\ref{tab:main_results}, using FFT+SigLIP2+Qwen features. Rows give the training dataset, columns the test dataset.}
\label{tab:transfer}
\resizebox{\columnwidth}{!}{%
\begin{tabular}{lccc}
\toprule
Train $\rightarrow$ Test & GenBuster & BigGAN & SD14 \\
\midrule
GenBuster & \textbf{0.899} & 0.995 & 0.874 \\
BigGAN    & 0.854 & \textbf{0.992} & 0.793 \\
SD14      & 0.899 & 0.995 & \textbf{0.874} \\
\midrule
Disagreement rule (no training) & 0.899 & 0.995 & 0.874 \\
\bottomrule
\end{tabular}
}
\end{table}

The GenBuster in-domain cell above (classifier accuracy 0.899) is byte-identical to the
disagreement rule's own accuracy on the same held-out subset (paired-bootstrap 95\% CI
on the difference: $[0.000, 0.000]$, 10{,}000 resamples): the learned classifier and the
parameter-free disagreement rule produce the same predictions. A previously reported
internal-split protocol, reproduced on current data as a robustness check, gives 0.900
(classifier) vs.\ 0.903 (disagreement rule), CI $[-0.0100, 0.0000]$ -- a small, likely
sampling-driven difference in the opposite direction, but one that is also not
distinguishable from zero. Both protocols therefore support the same conclusion: this
learned classifier provides no measurable improvement over the disagreement rule it was
trained to approximate, now with an explicit significance test rather than a bare point
comparison.


\subsubsection{Sensitivity to Semantic Probe Quality}

The disagreement signal assumes that both cheap detectors provide useful predictions.
When one detector falls to near-random accuracy, disagreement no longer identifies hard
cases reliably. Cross-family experiments show that semantic probes trained on one
generator family transfer poorly to distant families, with off-diagonal performance
remaining close to chance regardless of source domain.

This limits how the cascade should be deployed. Applying it to a new generator family
requires an in-domain semantic probe, or at least a probe that remains meaningfully above
chance on the target distribution. The disagreement rule adds no trainable parameters,
but it still depends on both cheap detectors that carry a usable signal.


\subsubsection{Cross-Family Deployment}
\label{sec:cross_family_deployment}

Cross-family deployment tests whether the cascade helps when the semantic probe is
applied outside its training family. A SigLIP2 probe trained on dataset A is evaluated
on dataset B without retraining. We compare the stale probe alone, an FFT-only detector
trained in-domain on the target dataset, and the disagreement cascade with Qwen as the
verifier (Table~\ref{tab:cross_family_deployment}). For reference, in-domain SigLIP2
performance is strong across all three datasets, so the off-diagonal tests isolate the
effect of cross-family degradation.

The cascade improves over the stale SigLIP2 probe in all six pairings, but that baseline
alone overstates the value of escalation. Against the in-domain FFT detector, four
pairings are significantly worse, one (SD~1.4$\to$GenBuster) is indistinguishable from
zero, and one (BigGAN$\to$GenBuster) is a small significant gain of 3.0 points
(95\% CI $[+0.005, +0.055]$).

These results illustrate \textbf{Case 2} from Section~III.D. The transferred semantic
probe is weak, so the cascade appears to recover accuracy relative to that stale probe.
However, many disagreement cases are inputs that the in-domain FFT detector already
classifies correctly. Escalating those cases to Qwen replaces correct FFT predictions
with a less reliable verifier prediction. Under this cross-family setting, a strong
in-domain artifact detector is therefore preferable to escalating disagreements from a
poorly transferred semantic probe.

This setting differs from the SD~1.4-to-SDXL/FLUX deployment-shift experiment in
Table~\ref{tab:e8_results}. Here, the semantic probe is transferred across generator
families and Qwen is used as the verifier. In Table~\ref{tab:e8_results}, the semantic
probe is trained on SD~1.4 and evaluated on newer diffusion models from the same lineage,
with GPT-5.5 as verifier. The positive cascade result in that setting is therefore a
different Case~3 configuration.

\begin{table}[h]
\centering
\caption{Cross-family deployment results. Cascade uses a cross-family SigLIP2 probe with Qwen escalation on disagreement, while FFT is trained in-domain on each evaluation dataset. $\Delta_{\text{stale}}$ denotes improvement over the stale SigLIP2 probe, and $\Delta_{\text{FFT}}$ denotes improvement over FFT-only. $^\dagger$ marks $\Delta_{\text{FFT}}$ cells whose paired-bootstrap 95\% CI includes zero (10,000 resamples per cell).}
\label{tab:cross_family_deployment}
\footnotesize
\setlength{\tabcolsep}{3pt}
\resizebox{\columnwidth}{!}{%
\begin{tabular}{llccccc}
\toprule
Train & Test & FFT & Stale SigLIP2 & Casc. & $\Delta_{\text{stale}}$ & $\Delta_{\text{FFT}}$ \\
\midrule
GenBuster & SD~1.4 & 0.775 & 0.512 & 0.636 & $+$0.124 & $-$0.140 \\
GenBuster & BigGAN & 0.997 & 0.530 & 0.776 & $+$0.246 & $-$0.221 \\
SD~1.4 & GenBuster & 0.727 & 0.502 & 0.702 & $+$0.200 & $-$0.025$^\dagger$ \\
SD~1.4 & BigGAN & 0.997 & 0.638 & 0.821 & $+$0.184 & $-$0.176 \\
BigGAN & GenBuster & 0.727 & 0.523 & 0.757 & $+$0.234 & $+$0.030 \\
BigGAN & SD~1.4 & 0.775 & 0.515 & 0.608 & $+$0.093 & $-$0.167 \\
\midrule
\multicolumn{2}{l}{Mean} & & & & $+$0.180 & $-$0.116 \\
\bottomrule
\end{tabular}
}
\end{table}

\begin{figure*}[t]
    \centering
    \includegraphics[width=0.7\textwidth]{plot_fig_a_semantic_degradation.png}
    \caption{Semantic probe accuracy under in-domain (solid) and deployment-shift evaluation (striped) on SDXL and FLUX.}
    \label{fig:semantic_degradation}
\end{figure*}

\subsubsection{Distant Generator Shift: BigGAN}
\label{sec:biggan_distant_shift}

SDXL and FLUX are diffusion models from the same broad lineage as SD~1.4, the dataset
used to train the stale FFT and semantic probes. To test a more distant shift, we apply
the SD~1.4-trained FFT and SigLIP2 probes zero-shot to GenImage BigGAN, a GAN-based
generator. This differs from the cross-family deployment experiment
(Section~V.D.6), where FFT is trained in-domain on the target dataset. Here, both cheap
detectors are stale, matching the SDXL/FLUX deployment-shift setup.

The BigGAN shift degrades both cheap detectors more sharply than SDXL or FLUX. On the
disagreement subset, however, stale SigLIP2 remains strong while stale FFT collapses.
This is an asymmetric Case~2 pattern rather than the joint degradation needed for
Case~3.

We evaluate GPT-5.5, Qwen2.5-VL-7B, and Qwen2.5-VL-72B (4-bit quantized) on a
stratified random sample of 1{,}500 disagreement cases. All three verifiers produce a 
negative cascade gain relative to the stale semantic probe (Table~\ref{tab:biggan_shift}).
None exceeds stale SigLIP2 on the routed subset, so escalation replaces reliable
semantic predictions with weaker verifier outputs. This matches the Case~2 failure mode
from the cross-family experiment (Section~V.D.6).

This negative result shows that the SDXL/FLUX finding does not extend automatically to
a distant generator family. The 72B model performs best among the tested verifiers yet
still falls far below the semantic probe, so verifier suitability cannot be inferred
from scale or general capability.

\begin{table}[h]
\centering
\caption{Verifier accuracy versus the best cheap detector on the BigGAN distant-shift disagreement subset ($n{=}1{,}500$).}
\label{tab:biggan_shift}
\footnotesize
\setlength{\tabcolsep}{3pt}
\begin{tabular}{l c c}
\toprule
Verifier & Acc. on Disagree & Best-cheap (Stale SigLIP2) \\
\midrule
GPT-5.5 & 0.582 & 0.907 \\
Qwen2.5-VL-7B & 0.515 & 0.907 \\
Qwen2.5-VL-72B & 0.623 & 0.907 \\
\bottomrule
\end{tabular}
\end{table}


\section{Discussion}

This section interprets the empirical results presented in Section V. We first examine
why artifact-semantic disagreement separates easy from hard inputs, then map the
three K-stage cascade cases from Section~III.D onto the empirical results and assess whether
learned routing improves on the disagreement signal, then characterize the role of the
semantic probe and VLM verifier under in-domain and deployment-shift conditions, and
finally discuss generator distribution shift, limitations, and practical deployment
guidelines.



\subsection{Artifact-Semantic Disagreement as a Routing Signal}

The complementarity arises from the signal types each detector uses. The FFT probe
captures spectral regularities introduced during synthesis, including upsampling grids
and DCT coefficient patterns. Semantic probes encode high-level scene content learned
from natural image-text distributions. These signal types are largely orthogonal.
Agreement indicates that these independent cues support the same label, while
disagreement exposes reduced reliability without calibrated confidence. The severity of
that reduction varies by dataset, so disagreement identifies candidates for escalation
without establishing that a verifier will improve them.

\subsection{The Three Cascade Cases in Practice}
\label{sec:three_cases}

The K-stage formulation (Section~III.D) predicts three qualitative cascade cases, and
the empirical results exhibit all three. \textbf{Case 1} occurs on BigGAN in-domain.
The artifact detector is already near ceiling, so the cascade adds latency with almost
no accuracy change (Table~\ref{tab:main_results}).

\textbf{Case 2} occurs when the cheap detectors are asymmetric. This describes the
in-domain GenBuster and SD~1.4 results (Section~V.B), and appears more clearly in the
cross-family deployment experiment (Table~\ref{tab:cross_family_deployment},
Section~V.D.6). A weak transferred probe makes the cascade appear useful against the
stale baseline, while comparison with the stronger cheap detector removes that apparent
gain. The BigGAN distant shift is another Case~2 instance because the semantic probe
remains stronger than every tested verifier on the routed subset.

\textbf{Case 3} occurs on the SD~1.4-to-SDXL/FLUX shift with GPT-5.5 as verifier
(Table~\ref{tab:e8_results}). Both cheap detectors degrade, and GPT-5.5 is stronger on
each routed subset. Replacing it with Qwen makes every gain negative, showing that the
case assignment depends on the full detector-verifier configuration.

Model scale alone does not determine the case. The clearest evidence is the
BigGAN distant-shift comparison (Table~\ref{tab:biggan_shift}), where Qwen2.5-VL-72B
outperforms GPT-5.5 yet remains far below the semantic probe. Both cheap detectors being
weak is insufficient for Case~3; the verifier must exceed the stronger cheap detector on
the routed subset, as Eq.~\eqref{eq:gain} requires.

The probe-degradation sweep (Table~\ref{tab:limited_data}) demonstrates the
Case~2-to-Case~3 transition within a controlled experiment. As the semantic probe
weakens, verifier escalation changes from harmful to useful. The absence of the same
transition on GenBuster confirms that relative routed-subset strength, rather than the
training fraction itself, determines the case.

The \textbf{breakdown condition} occurs when both cheap detectors share a blind spot.
Such errors remain invisible to disagreement-based routing by construction
(Table~\ref{tab:escalation}), and the agreement-path audit in
Section~\ref{sec:validation_gate} quantifies how often this occurs across every
evaluated dataset, not only GenBuster.

Figure~\ref{fig:case_mapping} summarizes the case assignment above across every
tested configuration in this paper, mapping Case 1/2/3/Breakdown onto each
dataset, verifier, and backbone combination discussed in Sections~V and~VI.

\begin{figure}[h]
    \centering
    \includegraphics[width=.95\columnwidth]{plot_fig_f_case_mapping.png}
    \caption{The Case 1/2/3/Breakdown framework (Section~III.D) mapped onto every tested configuration in this paper.}
    \label{fig:case_mapping}
\end{figure}

\subsection{Limitations of Learned Routing}

A learned escalation classifier with 11 detector-derived features and an explicit latency reward does
not provide a consistent improvement over binary disagreement. This suggests that the
dominant routing signal is already captured by whether the artifact and semantic
detectors agree, and that the learned policy mainly recovers the same decision boundary
rather than discovering a stronger routing rule (Table~\ref{tab:transfer}).


\subsection{Role of the Semantic Probe}

The role of the FFT detector in the cascade is diagnostic, not predictive. In agreement
cases, the FFT and semantic probe produce identical predictions by construction. Any
routing policy that escalates only disagreement cases to the VLM reduces to outputting
the semantic probe's prediction on agreement cases unchanged. The FFT detector
 does not contribute to independent prediction on these samples. It serves as a cheap signal
whose agreement with the semantic probe identifies easy inputs, routing only the
residual disagreement cases to the VLM.

This role makes the cascade most useful when the semantic probe becomes stale while the
verifier remains strong. Cross-family results also show that a parameter-free routing
rule still depends on useful cheap-detector predictions; a new generator family may
require an in-domain probe.

\subsection{VLM Domain Sensitivity}

Using the prompt in Section~IV.B with REAL/FAKE logit extraction, Qwen2.5-VL achieves
more reliable binary predictions when classification uses the REAL/FAKE logit gap rather
than parsed output tokens. Greedy decoding can bias the textual response toward
``REAL'' even when the underlying logits separate the classes. This is a methodological
consideration for VLM-based binary classification beyond this detector.

Even with logit extraction, Qwen remains weaker than the trained semantic probes
in-domain. Selective invocation improves on uniform Qwen use, though it still does not
surpass the semantic-only baseline. Routing can therefore improve how a weak verifier is
used without making it the best available model.

Qwen's routed-subset accuracy varies substantially across settings. In several
cross-family pairings it falls below the in-domain FFT baseline, explaining why
escalation does not reliably beat the stronger cheap detector. Verifier value is therefore
distribution-dependent.

\begin{figure}[h]
    \centering
    \includegraphics[width=.95\columnwidth]{plot_fig_b_cascade_gain_scatter.png}
    \caption{Cascade gain over the stale semantic probe versus stale semantic accuracy on SDXL and FLUX. Each point corresponds to one semantic backbone and deployment dataset.}
    \label{fig:cascade_gain}
\end{figure}

Figure~\ref{fig:cascade_gain} shows the same pattern for GPT-5.5 under SDXL and FLUX
deployment shift. Cascade gain generally increases as the stale semantic probe becomes
less accurate, though this relationship alone does not guarantee improvement. Verifier
strength should therefore be validated on the routed subset before enabling escalation.

\subsection{Generator Distribution Shift}

The contrast between CNNSpot and FFT suggests that artifact representation matters as
much as detector class under generator shift. CNNSpot's GAN-specific cues transfer
poorly to diffusion and video generators~\cite{Wang_2023_ICCV}, while the FFT probe's
spectral representation transfers more broadly. Artifact robustness alone still does
not determine cascade value because the routed verifier may remain weaker than the
semantic probe.




\subsection{Limitations}

One validity concern is the real-image source used in the SDXL/FLUX deployment-shift
evaluation. The real half of these benchmarks uses ImageNet val/nature photographs
rather than the bundled AIGIBench PASCAL VOC images (Section~IV.B). This means the real
images come from the same source as the real images used to train the SD~1.4 semantic
probes, while the fake SDXL/FLUX images remain out of distribution. We report results
on the original PASCAL VOC images and on an independently chosen real-image source
(COCO val2017) alongside this benchmark in Table~\ref{tab:real_image_sources}, rather
than asserting the substitution is innocuous based on class-wise analysis of the
substituted benchmark alone.

A reconciliation of this paper's own reported numbers uncovered a provenance bug in an
earlier draft's class-wise figures: the Limitations-section class-wise accuracy and
Table~\ref{tab:e8_results}'s overall cascade accuracy had been computed against two
different trained-probe snapshots saved roughly two hours apart, so the two sets of
numbers described two subtly different pipeline states rather than one consistent
result. We rebuilt both from a single, current, reproducible snapshot, with complete
(100\%) GPT-5.5 coverage of the disagreement subset via the OpenAI Batch API
(Section~\ref{sec:gpt55_audit}). Under this corrected, internally consistent accounting,
both real-class and fake-class recall \emph{increase} under cascade escalation, in all
four backbone/dataset cells: real recall rises from 95.1\%/91.4\% (SigLIP2/DINOv2) to
98.3\%/97.5\% on SDXL, and from 94.6\%/91.2\% to 98.4\%/97.5\% on FLUX; fake recall
rises from 77.7\%/78.9\% to 86.7\%/88.5\% on SDXL, and from 86.7\%/72.4\% to 89.7\%/81.1\%
on FLUX. This both corrects and strengthens an earlier draft's claim (which, under the
since-identified stale probe snapshot, reported real-class recall \emph{declining} by
4.0--4.3 points under escalation): under the current, reproducible data, the cascade's
accuracy gain is not confined to the fake class, and the real-image substitution does
not work against cascade accuracy as previously stated.

The primary evaluation sets contain 1{,}000 samples for GenBuster and 2{,}000 each
for SD~1.4 and BigGAN. Correlated cheap-detector errors affect 2.4\% of GenBuster
inputs, and 8.2--16.1\% of the SDXL/FLUX agreement path under deployment shift
(Section~\ref{sec:validation_gate}), showing that the routing signal's fundamental
limit is non-negligible and more pronounced under distribution shift than in-domain.
Larger sets would still strengthen generalization
claims. The learned escalation classifier is
trained and evaluated within the same benchmark distribution, so its generalization
to unseen generator types remains untested. Qwen2.5-VL is applied zero-shot without
fine-tuning. However, a fine-tuned VLM may recover the domain sensitivity failures observed on
SD~1.4 and BigGAN. The study uses frame-level detection without temporal modeling,
and incorporating motion and frame consistency could improve performance on
video-specific generators; the title and conclusions of this work should be read as
addressing single-frame forensic detection, not validated temporal video forensics. The semantic probe does not transfer across generator
families and deploying this cascade on a new family requires retraining the probe on
in-domain data. The disagreement rule itself has no trainable parameters and needs
no retraining. The SDXL/FLUX analysis excludes CLIP because the raw SD~1.4 features
required to reconstruct the stale probe were unavailable.

A further limitation concerns the artifact detector's near-ceiling in-domain accuracy
on BigGAN (99.7\%, Table~\ref{tab:main_results}). We found that BigGAN's fake images
are, without exception in a 300-image sample, natively exactly 128$\times$128 pixels,
while the corresponding real (ImageNet) images span over 300 distinct native
resolutions. A classifier given only an image's width and height, with no access to
pixel content, reaches 99.55\% accuracy on this same split -- within 0.15 points of the
trained FFT detector. Forcing the real class through the same native-resolution
downsampling the fake class structurally undergoes before the shared
256$\times$256-resize feature pipeline collapses the \emph{same trained} FFT detector's
accuracy from 99.7\% to 50.1\%, exact chance. A robustness check
(Table~\ref{tab:robustness}) shows the detector's BigGAN accuracy is also nearly
invariant to JPEG compression (99.7\%$\to$99.65\% at quality 50), consistent with a
coarse, compression-resistant resize signature rather than fine generative texture,
while accuracy collapses sharply under Gaussian blur and aggressive downscale-upscale
perturbations that directly alter spatial frequency content (99.7\%$\to$29.2\% at
blur $\sigma{=}2$; 99.7\%$\to$33.9\% at a 4$\times$ downscale-upscale round-trip),
including dropping below chance -- a pattern more consistent with a detector keyed to a
narrow, dataset-specific artifact than one tracking genuine generative texture. SD~1.4
shows a partial version of the same confound (resolution-only control 92.8\% vs.\ the
trained detector's 77.5\%; resolution-matching reduces it to 63.9\%, well above chance,
indicating a genuine signal persists alongside the confound); GenBuster, where both
classes share an identical native resolution, shows no such effect and serves as a
clean negative control (resolution-only control at exact chance, 50.0\%). This finding
has a direct consequence for Section~\ref{sec:cross_family_deployment}: two of the six
cross-family pairings there (GenBuster$\to$BigGAN and SD~1.4$\to$BigGAN) use an
in-domain-trained BigGAN FFT detector as the comparator that out-performs the cascade,
and that comparator's strength should now be read with this confound in mind rather
than as unqualified evidence of a strong, generalizable forensic detector.
Figure~\ref{fig:resolution_shortcut} summarizes this evidence visually.

Table~\ref{tab:robustness} extends the robustness check to all three cheap detectors.
FFT is the most fragile on every benchmark: blur at $\sigma{=}2$ and a 4$\times$ rescale
round trip push it to or below chance on SD~1.4 and BigGAN, and to 51--54\% on
GenBuster. The semantic probes degrade far less. SigLIP2 stays within 1 point of its
clean accuracy on GenBuster under every perturbation and loses at most 3 points on
SD~1.4, though it drops to 81--83\% on BigGAN under the strongest blur and rescale.
JPEG compression costs SigLIP2 at most 0.8 points and DINOv2 up to 4 points, and hits
FFT hardest on SD~1.4 (77.5\%$\to$67.1\% at quality 50). Because disagreement routing depends on the artifact detector, these
perturbations would raise the disagreement rate and move more inputs to the verifier.

\begin{table*}[t]
\centering
\footnotesize
\caption{Cheap-detector accuracy under post-processing on the in-domain evaluation sets: JPEG quality $q$, Gaussian blur $\sigma$, and a down/up-scale round trip by factor $s$.}
\label{tab:robustness}
\begin{tabular}{llcccccccc}
\toprule
Dataset & Detector & Clean & $q{=}95$ & $q{=}75$ & $q{=}50$ & $\sigma{=}1$ & $\sigma{=}2$ & $s{=}0.5$ & $s{=}0.25$ \\
\midrule
GenBuster & FFT & 0.727 & 0.724 & 0.722 & 0.725 & 0.660 & 0.511 & 0.710 & 0.537 \\
GenBuster & SigLIP2 & 0.966 & 0.967 & 0.966 & 0.963 & 0.967 & 0.968 & 0.966 & 0.966 \\
GenBuster & DINOv2 & 0.953 & 0.953 & 0.946 & 0.933 & 0.954 & 0.951 & 0.951 & 0.953 \\
\midrule
SD~1.4 & FFT & 0.775 & 0.795 & 0.711 & 0.671 & 0.683 & 0.493 & 0.647 & 0.491 \\
SD~1.4 & SigLIP2 & 0.926 & 0.926 & 0.936 & 0.937 & 0.917 & 0.897 & 0.917 & 0.907 \\
SD~1.4 & DINOv2 & 0.877 & 0.866 & 0.835 & 0.837 & 0.849 & 0.819 & 0.859 & 0.831 \\
\midrule
BigGAN & FFT & 0.997 & 0.997 & 0.996 & 0.997 & 0.883 & 0.292 & 0.952 & 0.339 \\
BigGAN & SigLIP2 & 0.992 & 0.988 & 0.984 & 0.988 & 0.975 & 0.807 & 0.977 & 0.833 \\
BigGAN & DINOv2 & 0.984 & 0.975 & 0.957 & 0.945 & 0.974 & 0.924 & 0.963 & 0.902 \\
\bottomrule
\end{tabular}
\end{table*}

\begin{figure}[h]
    \centering
    \includegraphics[width=.95\columnwidth]{plot_fig_e_resolution_shortcut.png}
    \caption{A width/height-only classifier (no pixel content) matches BigGAN's trained FFT detector; resolution-matching collapses the trained detector to chance. SD~1.4 shows a partial confound; GenBuster (no resolution difference between classes) is a clean negative control.}
    \label{fig:resolution_shortcut}
\end{figure}

\subsection{Practical Implications}

Deployment should begin with the strongest cheap detector and treat disagreement as a
candidate escalation trigger. If a semantic probe is accurate in-domain, direct probe
inference is preferable. Under shift, escalation should be enabled only after the chosen
verifier demonstrates an advantage on a labeled sample from the routed subset. This
policy separates the identification of difficult inputs from the question of whether a
particular verifier can resolve them.

\subsection{A Practical Fail-Safe: the Verifier Validation Gate}
\label{sec:validation_gate}

A deployed system may not know whether an incoming distribution is in-domain, a benign
generator update, or a larger cross-family shift. A natural first option is to use the
disagreement rate as a shift signal. This is insufficient: the cross-family and
SDXL/FLUX ranges overlap despite opposite escalation outcomes, while the BigGAN distant
shift remains close to the in-domain rate even though every tested verifier hurts
performance there.

We therefore propose a \textbf{verifier validation gate}. Before allowing escalation for
a verifier $V$ on a new target distribution, evaluate $V$ on a small labeled calibration
sample drawn from the disagreement set, and freeze the resulting decision before
applying it to new, unlabeled inputs.

\textbf{Decision rule.} We specify two versions of the gate, in pseudocode, to make
their relationship precise rather than implicit (this precision is necessary: an
earlier version of this section stated the two were different in a way a careful
reader showed was incorrect -- see below).

\begin{verbatim}
# Binary gate: escalate-or-fallback
def binary_gate(acc1, acc2, acc3, ci_lower):
    if acc3 > max(acc1, acc2):           # naive (point-estimate)
    # or: if ci_lower(acc3 - max(acc1,acc2)) > 0:  # significance-aware
        return ESCALATE_TO_VERIFIER
    else:
        return USE_STRONGER_CHEAP_DETECTOR  # i.e. argmax(acc1, acc2)

# Three-way selector
def three_way_selector(acc1, acc2, acc3):
    return argmax([acc1, acc2, acc3])    # f1, f2, or f3
\end{verbatim}

As decision procedures, these two rules select the identical action in every
case: the binary gate's fallback branch is, by its own definition,
$\arg\max(A_1, A_2)$, so "escalate if $A_3$ beats $\max(A_1,A_2)$, else use
whichever of $f_1/f_2$ is larger" is exactly $\arg\max(A_1, A_2, A_3)$ restricted
to a binary escalate/fallback framing. A reviewer correctly pointed out that the
previous version of this paper's claim that the binary rule "cannot represent
artifact-detector selection" conflicts with this definition; it does. We retract
that claim. What the two rules differ in is not the action selected but what the
accompanying evaluation reported: our original calibration-sensitivity analysis
(below) tracked only the binary rule's escalate/fallback Boolean against the
oracle, not which of $f_1$/$f_2$ it selected on fallback, while the three-way
analysis tracked the full three-valued identity. This is a difference in what was
measured, not in what the gate does, and it fully explains the one place the two
rules' reported numbers diverge (at $N{=}10$; see below).

We use a prespecified, uncertainty-aware version of the escalate condition rather
than the bare point estimate: escalate only if the paired-bootstrap 95\% confidence
interval (10{,}000 resamples) on $(A_{V} - \max(A_1, A_2))$, computed on the
calibration sample, lies entirely above zero. A small positive point-estimate gap
that could plausibly be sampling noise no longer triggers escalation.

The calibration/evaluation set for this gate spans the same 35 configurations used
throughout Section~V's baseline comparison: the verifier-on-disagreement-subset
cells from Tables~\ref{tab:verifier_comparison}, \ref{tab:e8_results},
\ref{tab:cross_family_deployment}, \ref{tab:limited_data}, and
\ref{tab:biggan_shift}.

\textbf{Agreement-path audit.} Calibrating only on the disagreement subset cannot by
itself establish the accuracy of the agreement path, where no verifier is ever
consulted. We separately audited agreement-path accuracy and the rate of correlated
cheap-detector errors (both detectors agree and are both wrong -- the one failure
mode disagreement-based routing structurally cannot catch) across all evaluated
datasets. In-domain agreement-path accuracy is high and correlated-error rates are
low (GenBuster 96.8\% / 3.2\%, SD~1.4 97.8\% / 2.2\%, BigGAN 99.8\% / 0.2\%, accuracy
/ correlated-error rate respectively). Under SDXL/FLUX deployment shift, agreement-path
accuracy drops materially and correlated-error rates rise 4--8$\times$ (SDXL SigLIP2
88.4\% / 11.6\%, SDXL DINOv2 88.7\% / 11.3\%, FLUX SigLIP2 91.8\% / 8.2\%, FLUX DINOv2
83.9\% / 16.1\%): a meaningful share of agreement-path predictions under deployment
shift are errors the gate's calibration, restricted to the disagreement subset, cannot
see. This is the empirical shape of the breakdown condition from Section~III.D under
deployment shift specifically, not a universal property of the agreement path.

\subsubsection{Independent (Held-Out) Gate Evaluation}
\label{sec:gate_holdout}

Validating the gate's decision against an oracle computed from the \emph{same} pool
used for calibration establishes resampling stability, not generalization to unseen
inputs. We therefore also evaluated the gate under genuine calibration/test
separation. For each of the 35 configurations and each calibration size
$N \in \{50, 100, 200\}$, we drew 250 random splits of the disagreement subset into
disjoint calibration and test sets of $N$ samples each, plus 250 splits into two
halves of the full pool. When a pool holds fewer than $2N$ samples, $N$ is capped at
half the pool: GenBuster/GPT-5.5 ($n{=}52$) is capped at 26 for every $N$, and
GenBuster/Qwen2.5-VL-7B ($n{=}259$) at 129 for $N{=}200$. On each split we froze the
gate's decision on the calibration set, under both the naive point-estimate rule and
the significance-aware rule above (500 bootstrap resamples per draw in this
simulation), and applied it unchanged to the test set. We record agreement with the
test set's own oracle decision, regret against the best fixed detector chosen with
test-set knowledge, harmful escalations (the frozen decision escalates but the
verifier is below the best cheap detector on the test set), and missed beneficial
escalations (the verifier is strictly best on the test set but the frozen decision
falls back). Three of the 35 configurations have disagreement pools under 40 samples
(SD~1.4/GPT-5.5, $n{=}23$; BigGAN/Qwen2.5-VL-7B and BigGAN/GPT-5.5, both $n{=}14$) and
cannot support a meaningful split; we report the remaining 32.

Under this protocol (Table~\ref{tab:gate_holdout}), the significance-aware rule
reaches $\geq$90\% agreement with the test oracle in 18/32 configurations at
$N{=}50$, 21/32 at $N{=}100$ and 26/32 at $N{=}200$ (naive rule: 17/32, 21/32,
26/32). These counts are below the same-pool figures of 22/35, 24/35 and 28/35 in
Table~\ref{tab:sensitivity}, confirming that the earlier numbers measured resampling
stability rather than held-out generalization. We revise the practical recommendation
accordingly. $N{=}100$ gives reliable decisions in only about two-thirds of eligible
configurations, while $N{=}200$ reaches about four-fifths, so practitioners should
budget closer to $N{\approx}200$ or use the full calibration pool, and should still expect some close-margin configurations to
stay unreliable.

The two rules make a different trade, not a better or worse one. The
significance-aware rule cuts harmful escalations by roughly an order of magnitude
(6.6\% vs.\ 0.2\% at $N{=}50$). It pays for this by escalating less: in 11--17\% of
draws it falls back although the verifier was best on the test set, against 2--8\%
for the naive rule, and its mean regret is higher at every $N$. This is also why it
matches the test oracle less often than the naive rule with the half pool (26 vs.\
28). With that much calibration data the naive rule's
point estimate is usually right, while the confidence interval still excludes zero
only when the verifier's margin is large. The significance-aware rule suits
deployments where an unwarranted escalation is costly. The naive rule suits those
where a missed verifier gain is.

The lower panel of Table~\ref{tab:gate_holdout} shows how often the
significance-aware gate enables escalation on the four SDXL/FLUX GPT-5.5
configurations behind the paper's positive result. Three of them escalate in at least
85\% of draws at every $N$. FLUX/SigLIP2, where the verifier's margin over the stale
probe is smallest, escalates in only 38\% of draws at $N{=}50$ and 64\% at $N{=}100$,
rising to 94\% at $N{=}200$ and 100\% with the half pool. With 100 or fewer labeled
calibration samples, the gate would therefore forgo the FLUX/SigLIP2 gain in a third
or more of deployments.

\subsubsection{Calibration Sample-Size Sensitivity (Same-Pool Oracle)}

The sensitivity sweep below retains the earlier same-pool-oracle protocol (calibrating
against, and evaluating against, the same full disagreement subset) alongside the
held-out results above, since it isolates resampling stability as a separate question
from generalization and was the basis of our original (now revised) sample-size
recommendation. For all 35 configurations, we bootstrap-subsampled each disagreement subset
at $N \in \{10, 20, 50, 100, 200\}$ (1{,}000 resamples per configuration-$N$
pair) and measured how often the gate's decision from the subsample matches
the oracle decision computed from the full pool. Table~\ref{tab:sensitivity} summarizes the results.

\FloatBarrier
\begin{table}[t]
\centering
\small
\caption{Same-pool-oracle sensitivity: fraction of the 35 configurations reaching $\geq$90\% oracle-decision agreement at each calibration sample size $N$, for the binary escalate/fallback gate and the three-way $f_1$/$f_2$/$f_3$ selector (Section~\ref{sec:threeway}). Measures resampling stability against the SAME pool used to define the oracle, not generalization; see Table~\ref{tab:gate_holdout} for the held-out counterpart. \textit{Note: for configurations whose disagreement subset contains fewer than 200 unique samples, the $N{=}100$/$N{=}200$ columns are bootstrap resamples with replacement from that smaller pool, not a stand-in for 100 or 200 independent human-labeled examples.}}
\label{tab:sensitivity}
\begin{tabular}{lccccc}
\toprule
Decision rule & $N{=}10$ & $N{=}20$ & $N{=}50$ & $N{=}100$ & $N{=}200$ \\
\midrule
Binary gate         & 14/35 & 19/35 & 22/35 & 24/35 & 28/35 \\
Three-way selector  & 13/35 & 19/35 & 22/35 & 24/35 & 28/35 \\
\bottomrule
\end{tabular}
\end{table}
\FloatBarrier

\FloatBarrier
\begin{table}[t]
\centering
\small
\caption{Independent (held-out) gate evaluation over 32 eligible configurations (3 of 35 excluded, disagreement pool $<$40), 250 disjoint calibration/test draws per configuration and $N$. Upper panel: configurations reaching $\geq$90\% agreement with the test-set oracle, then the mean rates of harmful escalation (escalated, verifier below best cheap detector on test) and missed beneficial escalation (fell back, verifier strictly best on test), and mean regret. Lower panel: share of draws in which the significance-aware gate enables escalation on the SDXL/FLUX GPT-5.5 configurations. $N$ is capped at half the pool for GenBuster/GPT-5.5 (26 at every $N$) and GenBuster/Qwen2.5-VL-7B (129 at $N{=}200$).}
\label{tab:gate_holdout}
\resizebox{\columnwidth}{!}{%
\begin{tabular}{lcccc}
\toprule
Decision rule & $N{=}50$ & $N{=}100$ & $N{=}200$ & Half-pool \\
\midrule
Naive (point-estimate)         & 17/32 & 21/32 & 26/32 & 28/32 \\
\ \ harmful-esc.\ rate         & 6.6\% & 5.5\% & 4.0\% & 2.2\% \\
\ \ missed-benefit rate        & 8.1\% & 6.4\% & 3.5\% & 2.0\% \\
\ \ mean regret (points)       & 1.39 & 0.76 & 0.35 & 0.18 \\
Significance-aware             & 18/32 & 21/32 & 26/32 & 26/32 \\
\ \ harmful-esc.\ rate         & 0.2\% & 0.4\% & 0.5\% & 0.5\% \\
\ \ missed-benefit rate        & 17.1\% & 15.1\% & 12.5\% & 11.3\% \\
\ \ mean regret (points)       & 2.00 & 1.27 & 0.78 & 0.62 \\
Rules choose differently       & 16.5\% & 14.7\% & 12.8\% & 11.3\% \\
\midrule
\multicolumn{5}{l}{Significance-aware gate escalates (share of draws)} \\
\ \ SDXL/SigLIP2               & 96.0\% & 100.0\% & 100.0\% & 100.0\% \\
\ \ SDXL/DINOv2                & 98.4\% & 100.0\% & 100.0\% & 100.0\% \\
\ \ FLUX/SigLIP2               & 37.6\% & 64.0\% & 93.6\% & 100.0\% \\
\ \ FLUX/DINOv2                & 86.8\% & 99.2\% & 100.0\% & 100.0\% \\
\bottomrule
\end{tabular}
}
\end{table}
\FloatBarrier

\FloatBarrier
\begin{table}[h]
\centering
\small
\caption{Oracle three-way detector selection across 35 configurations.}
\label{tab:threeway}
\begin{tabular}{lc}
\toprule
Selected detector & Configurations \\
\midrule
Semantic probe ($f_2$) & 18 / 35 \\
Verifier ($f_3$)       &  9 / 35 \\
Artifact detector ($f_1$) & 8 / 35 \\
\bottomrule
\end{tabular}
\end{table}
\FloatBarrier

Under this same-pool-oracle protocol, both decision rules reach reliable
($\geq$90\%) agreement in the majority (24/35) of tested configurations at
$N{\geq}100$. As established in Section~\ref{sec:gate_holdout}, this same-pool
figure should not be read as a generalization guarantee; our practical
recommendation is based on the held-out evaluation there, not on this sweep alone.
For configurations whose
disagreement subset contains fewer than 200 unique samples, the $N=100$/$N=200$
draws above are bootstrap resamples with replacement from that smaller pool,
not a stand-in for 100 or 200 independent human-labeled examples. A subset of
close-margin configurations, where verifier and best-cheap-detector accuracy
are within a few points of each other, remain noisy even at $N=200$.
Practitioners deploying the gate on a configuration suspected to be close-margin
should use a larger calibration batch or treat the recommendation as provisional.

The three-way selector's agreement rate matches the binary gate's at every tested
$N$ except $N=10$, where it trails by one configuration (14/35 vs.\ 13/35). We traced
this to a single configuration (SD~1.4/SigLIP2 at 100\% training fraction, Qwen2.5-VL-7B
verifier; oracle accuracies $A_1{=}0.216$, $A_2{=}0.784$, $A_3{=}0.587$, with
$A_1 + A_2 = 1$ since every disagreement-subset item is, by construction, one the two
cheap detectors predict oppositely). At $N{=}10$, the resampled accuracy of a cheap
detector can only take 11 discrete values ($k/10$), so the resampled $\hat A_1$ can tie
or even exceed $\hat A_2$ although $A_1$ is far below $A_2$. A tie occurs in 29 of the
1{,}000 resamples, close to the $\mathrm{Binomial}(10, 0.216).\mathrm{pmf}(5) \approx
3.5\%$ expectation, and $\hat A_1 > \hat A_2$ in 14 more. \texttt{argmax} resolves ties
to the first-listed detector ($f_1$, the artifact detector). Which cheap detector is
chosen does not change the binary rule's escalate/fallback outcome, but it does change
the three-way rule's full-identity agreement check. Of the 1{,}000 resamples, 89 have
$\hat A_3 > \max(\hat A_1, \hat A_2)$, a mismatch both rules share. A further 27 select
$f_1$ where the oracle selects $f_2$, a mismatch only the three-way check counts:
$89 + 27 = 116 = 1{,}000 - 884$, which closes the arithmetic exactly against the
observed 91.1\% vs.\ 88.4\% agreement rates on this configuration. This
single-configuration resampling artifact, not a difference in the gate's underlying
decision, accounts for the entire $N{=}10$ discrepancy. It vanishes by $N{=}20$, where
the resampled accuracies are fine-grained enough that such reversals are rare.

\subsubsection{Extending the Gate to Three-Way Selection}
\label{sec:threeway}

The binary gate above chooses between escalating to the verifier and
falling back to the stronger cheap detector. We generalize this to a
three-way selection rule: calibrate the accuracy of $f_1$ (artifact),
$f_2$ (semantic), and $f_3$ (verifier) on the same labeled calibration
sample, and select the highest-performing detector for subsequent inputs
from that deployment's disagreement subset. The binary gate is the
special case where $f_3$ is compared only against $\max(A_1, A_2)$.

The oracle analysis shows that no single detector dominates across
deployment conditions (Table~\ref{tab:threeway}). In particular, the
artifact detector is preferred in five of the six cross-family pairings
(Table~\ref{tab:cross_family_deployment}). As established above, the binary gate's
fallback branch already selects this detector whenever it is the stronger of the two
cheap detectors, so the binary and three-way rules make the identical choice here;
the three-way framing's practical value is reporting that choice as an explicit,
separately-calibrated selection among all three detectors, rather than leaving it as
an implicit byproduct of a two-outcome escalate/fallback label.


The oracle result above uses the full disagreement subset as calibration.
With a realistic labeled batch, the three-way selector closely tracks the
binary gate's calibration reliability across the full sensitivity sweep
(Table~\ref{tab:sensitivity}). The selector reduces to the binary gate's
recommendation whenever the artifact detector is not competitive. Its practical
value is selecting among all three detectors directly, using the same calibration
sample as the binary gate.

\subsection{Cost Accounting and Cascade vs.\ Uniform Escalation}
\label{sec:cost_accounting}

Throughout this paper we report wall-clock latency and, where measured, GPT-5.5
token counts, but a direct cost-accuracy comparison of the cascade against both a
standalone cheap detector and a uniform (always-escalate) GPT-5.5 policy had not
been possible: no prior evaluation in this line of work included a full-coverage
GPT-5.5 run on an entire benchmark, only disagreement-subset-restricted calls. We
closed this gap for the two SDXL/FLUX configurations supporting this paper's
positive escalation result (Table~\ref{tab:e8_results}), using the OpenAI Batch API
(Batch-tier pricing of \$2.50/\$15.00 per million input/output tokens for gpt-5.5 at
standard context length~\cite{openai_pricing}, half the synchronous-API price,
appropriate here since none of this evaluation is latency-sensitive) to run GPT-5.5 on
every image in both benchmarks (8{,}000 calls, 100\% successful on first pass). A
GPT-5.5 call used on average 783 input and 96 output tokens on SDXL and 783 input and
177 output tokens on FLUX. In the in-domain disagreement-subset runs the averages were
1{,}271/233 (GenBuster), 339/211 (SD~1.4) and 192/184 (BigGAN) input/output tokens.

Table~\ref{tab:cost_accuracy_tradeoff} and Figure~\ref{fig:cost_accuracy_tradeoff}
report the resulting three-point comparison. Uniform escalation reaches 99.6\%
accuracy on SDXL and 97.1\% on FLUX, at \$13.61 and \$18.47 respectively for the full
4{,}000-image benchmark. The disagreement-routed cascade reaches 92.5--93.0\% (SDXL)
and 89.3--94.0\% (FLUX) at 42--51\% of that cost, escalating 36--42\% of inputs. The cost
share exceeds the escalation share because disagreement images use more tokens per call
than the average image (for example, 965 vs.\ 783 input tokens on SDXL/SigLIP2), so the
cascade's cost is computed from the tokens of the images it actually escalates. The
cascade therefore does not recover the full accuracy achievable by uniform escalation
-- 3.0 to 7.8 percentage points remain on the table across the four
backbone/dataset cells -- in exchange for a substantial, but not complete, cost
reduction. We report this explicitly rather than let the cascade's cost savings imply
it matches uniform escalation's accuracy: the appropriate reading is that selective
escalation is a genuine, quantifiable cost-accuracy tradeoff point, not a free
improvement over always escalating.

\begin{table}[h]
\centering
\footnotesize
\caption{Cheap-only, cascade, and uniform GPT-5.5 accuracy and cost on SDXL/FLUX (full 4{,}000-image benchmark, Batch-API pricing applied to per-call token counts). Cost share is cascade cost as a share of uniform cost.}
\label{tab:cost_accuracy_tradeoff}
\resizebox{\columnwidth}{!}{%
\begin{tabular}{llccccc}
\toprule
Dataset & Sem. & Cheap-only & Cascade (\$) & Uniform (\$) & Gap & Cost share \\
\midrule
SDXL & SigLIP2 & 0.864 & 0.925 (\$5.72) & 0.996 (\$13.61) & 0.071 & 42\% \\
SDXL & DINOv2  & 0.851 & 0.930 (\$6.01) & 0.996 (\$13.61) & 0.067 & 44\% \\
FLUX & SigLIP2 & 0.906 & 0.940 (\$9.40) & 0.971 (\$18.47) & 0.030 & 51\% \\
FLUX & DINOv2  & 0.818 & 0.893 (\$9.25) & 0.971 (\$18.47) & 0.078 & 50\% \\
\bottomrule
\end{tabular}
}
\end{table}

\begin{figure}[h]
    \centering
    \includegraphics[width=.95\columnwidth]{plot_fig_d_cost_accuracy_tradeoff.png}
    \caption{Cost vs.\ accuracy for cheap-only, cascade, and uniform GPT-5.5 escalation on SDXL and FLUX. The cascade recovers most, but not all, of the accuracy gap to uniform escalation at 42--51\% of its cost.}
    \label{fig:cost_accuracy_tradeoff}
\end{figure}

This cost accounting is necessarily incomplete in one respect: it covers GPT-5.5's
API cost only (Qwen2.5-VL runs locally and has no comparable \$ cost, reported
instead as latency throughout). All dollar figures use token counts measured
per-call and are clearly separated from the hardcoded, non-measured cheap-detector
latency constants (FFT 2.5~ms, CLIP/SigLIP2/DINOv2 10.4~ms, CNNSpot 45.8~ms,
Qwen always-escalate 190.2/83.2/70.6~ms for GenBuster/SD~1.4/BigGAN) used elsewhere
in this paper's cost comparisons, which remain fixed hardware measurements rather
than figures re-derived for this analysis.

\section{Conclusion}

This work frames AI-generated media detection as an adaptive model selection problem. A two-stage cheap pipeline routes inputs through an FFT artifact detector and a semantic probe and escalates to a frontier VLM only when the two detectors disagree. Three research questions were evaluated with this design.

For RQ1, artifact-semantic disagreement consistently identifies a subset with reduced
cheap-detector reliability, though the severity varies by dataset.

For RQ2, the disagreement rule does not require parameters or calibration. A learned
escalation classifier using the same detector outputs does not provide a consistent advantage, which
supports binary disagreement as the dominant routing signal in this setting. A joint
artifact-and-semantic confidence-routing baseline with held-out thresholds qualifies
this. The disagreement rule is more accurate on SD~1.4 and BigGAN in-domain, but under
the SDXL/FLUX shift the calibrated confidence rule is tied or better in every
configuration. Disagreement is a calibration-free routing signal that performs close to
a calibrated confidence rule, not a more accurate one.

For RQ3, disagreement-based escalation reduces verifier cost, but its accuracy benefit is
conditional. It recovers accuracy under the SDXL and FLUX shifts when paired with a
strong verifier, yet fails under cross-family transfer and the distant BigGAN shift.
Escalation should be enabled only when calibration confirms that the verifier exceeds
both cheap detectors on the routed subset.

The effective deployment of heterogeneous detection pipelines depends on the use of cheap detectors to identify when expensive verification is needed rather than increasing the complexity of the individual detector. Future work will explore temporal modeling for video detection, fine-tuned VLM verifiers for domain-specific forensic cues, and cross-family semantic probe training to remove the per-family retraining requirement.


\bibliographystyle{IEEEtran}
\bibliography{references}
\vskip -.5\baselineskip plus -1fil
\begin{IEEEbiography}[{\includegraphics[width=1in,height=1.25in,clip,keepaspectratio]{ansuman_nayak2.png}}]{Ansuman Nayak}
is a Software Engineer at Bank of America, where he focuses on building scalable automation tools for deployment pipelines, data workflows, and infrastructure operations. He earned his Bachelor's degree in Computer Science from Virginia Tech and is currently pursuing a Master's degree in Computer Science at Georgia Institute of Technology, specializing in Artificial Intelligence. His interests center on Agentic AI and LLM optimization, with a focus on workflow automation designed to operate efficiently under resource and cost constraints.
\end{IEEEbiography}
\vspace{-0.5\baselineskip}
\begin{IEEEbiography}[{\includegraphics[width=1in,height=1.25in,clip,keepaspectratio]{vijay.png}}]{Vijay K. Madisetti}
holds the position of Professor of Cybersecurity and Privacy (SCP) at Georgia Tech.
He earned his PhD in Electrical Engineering and Computer Sciences from the University of California at Berkeley
and is recognized as a Fellow of the IEEE. Additionally, he has been honored with the Terman Medal
by the American Society of Engineering Education (ASEE). Professor Madisetti has authored several widely
referenced textbooks on topics including cloud computing, data analytics, blockchain, and microservices.
\end{IEEEbiography}

\EOD

\end{document}